\documentclass[
  doc,
  floatsintext,
  longtable,
  a4paper,
  nolmodern,
  notxfonts,
  notimes,
  12pt,
  colorlinks=true,linkcolor=blue,citecolor=blue,urlcolor=blue]{apa7}

\usepackage{amsmath}
\usepackage{amssymb}

\geometry{inner=1in, outer=1in}
\fancyhfoffset[LE,RO]{0cm}


\usepackage[bidi=default]{babel}
\babelprovide[main,import]{spanish}


% get rid of language-specific shorthands (see #6817):
\let\LanguageShortHands\languageshorthands
\def\languageshorthands#1{}

\RequirePackage{longtable}
\RequirePackage{threeparttablex}

\makeatletter
\renewcommand{\paragraph}{\@startsection{paragraph}{4}{\parindent}%
	{0\baselineskip \@plus 0.2ex \@minus 0.2ex}%
	{-.5em}%
	{\normalfont\normalsize\bfseries\typesectitle}}

\renewcommand{\subparagraph}[1]{\@startsection{subparagraph}{5}{0.5em}%
	{0\baselineskip \@plus 0.2ex \@minus 0.2ex}%
	{-\z@\relax}%
	{\normalfont\normalsize\bfseries\itshape\hspace{\parindent}{#1}\textit{\addperi}}{\relax}}
\makeatother




\usepackage{longtable, booktabs, multirow, multicol, colortbl, hhline, caption, array, float, xpatch}
\usepackage{subcaption}


\renewcommand\thesubfigure{\Alph{subfigure}}
\setcounter{topnumber}{2}
\setcounter{bottomnumber}{2}
\setcounter{totalnumber}{4}
\renewcommand{\topfraction}{0.85}
\renewcommand{\bottomfraction}{0.85}
\renewcommand{\textfraction}{0.15}
\renewcommand{\floatpagefraction}{0.7}

\usepackage{tcolorbox}
\tcbuselibrary{listings,theorems, breakable, skins}
\usepackage{fontawesome5}

\definecolor{quarto-callout-color}{HTML}{909090}
\definecolor{quarto-callout-note-color}{HTML}{0758E5}
\definecolor{quarto-callout-important-color}{HTML}{CC1914}
\definecolor{quarto-callout-warning-color}{HTML}{EB9113}
\definecolor{quarto-callout-tip-color}{HTML}{00A047}
\definecolor{quarto-callout-caution-color}{HTML}{FC5300}
\definecolor{quarto-callout-color-frame}{HTML}{ACACAC}
\definecolor{quarto-callout-note-color-frame}{HTML}{4582EC}
\definecolor{quarto-callout-important-color-frame}{HTML}{D9534F}
\definecolor{quarto-callout-warning-color-frame}{HTML}{F0AD4E}
\definecolor{quarto-callout-tip-color-frame}{HTML}{02B875}
\definecolor{quarto-callout-caution-color-frame}{HTML}{FD7E14}

%\newlength\Oldarrayrulewidth
%\newlength\Oldtabcolsep


\usepackage{hyperref}



\usepackage{color}
\usepackage{fancyvrb}
\newcommand{\VerbBar}{|}
\newcommand{\VERB}{\Verb[commandchars=\\\{\}]}
\DefineVerbatimEnvironment{Highlighting}{Verbatim}{commandchars=\\\{\}}
% Add ',fontsize=\small' for more characters per line
\usepackage{framed}
\definecolor{shadecolor}{RGB}{241,243,245}
\newenvironment{Shaded}{\begin{snugshade}}{\end{snugshade}}
\newcommand{\AlertTok}[1]{\textcolor[rgb]{0.68,0.00,0.00}{#1}}
\newcommand{\AnnotationTok}[1]{\textcolor[rgb]{0.37,0.37,0.37}{#1}}
\newcommand{\AttributeTok}[1]{\textcolor[rgb]{0.40,0.45,0.13}{#1}}
\newcommand{\BaseNTok}[1]{\textcolor[rgb]{0.68,0.00,0.00}{#1}}
\newcommand{\BuiltInTok}[1]{\textcolor[rgb]{0.00,0.23,0.31}{#1}}
\newcommand{\CharTok}[1]{\textcolor[rgb]{0.13,0.47,0.30}{#1}}
\newcommand{\CommentTok}[1]{\textcolor[rgb]{0.37,0.37,0.37}{#1}}
\newcommand{\CommentVarTok}[1]{\textcolor[rgb]{0.37,0.37,0.37}{\textit{#1}}}
\newcommand{\ConstantTok}[1]{\textcolor[rgb]{0.56,0.35,0.01}{#1}}
\newcommand{\ControlFlowTok}[1]{\textcolor[rgb]{0.00,0.23,0.31}{\textbf{#1}}}
\newcommand{\DataTypeTok}[1]{\textcolor[rgb]{0.68,0.00,0.00}{#1}}
\newcommand{\DecValTok}[1]{\textcolor[rgb]{0.68,0.00,0.00}{#1}}
\newcommand{\DocumentationTok}[1]{\textcolor[rgb]{0.37,0.37,0.37}{\textit{#1}}}
\newcommand{\ErrorTok}[1]{\textcolor[rgb]{0.68,0.00,0.00}{#1}}
\newcommand{\ExtensionTok}[1]{\textcolor[rgb]{0.00,0.23,0.31}{#1}}
\newcommand{\FloatTok}[1]{\textcolor[rgb]{0.68,0.00,0.00}{#1}}
\newcommand{\FunctionTok}[1]{\textcolor[rgb]{0.28,0.35,0.67}{#1}}
\newcommand{\ImportTok}[1]{\textcolor[rgb]{0.00,0.46,0.62}{#1}}
\newcommand{\InformationTok}[1]{\textcolor[rgb]{0.37,0.37,0.37}{#1}}
\newcommand{\KeywordTok}[1]{\textcolor[rgb]{0.00,0.23,0.31}{\textbf{#1}}}
\newcommand{\NormalTok}[1]{\textcolor[rgb]{0.00,0.23,0.31}{#1}}
\newcommand{\OperatorTok}[1]{\textcolor[rgb]{0.37,0.37,0.37}{#1}}
\newcommand{\OtherTok}[1]{\textcolor[rgb]{0.00,0.23,0.31}{#1}}
\newcommand{\PreprocessorTok}[1]{\textcolor[rgb]{0.68,0.00,0.00}{#1}}
\newcommand{\RegionMarkerTok}[1]{\textcolor[rgb]{0.00,0.23,0.31}{#1}}
\newcommand{\SpecialCharTok}[1]{\textcolor[rgb]{0.37,0.37,0.37}{#1}}
\newcommand{\SpecialStringTok}[1]{\textcolor[rgb]{0.13,0.47,0.30}{#1}}
\newcommand{\StringTok}[1]{\textcolor[rgb]{0.13,0.47,0.30}{#1}}
\newcommand{\VariableTok}[1]{\textcolor[rgb]{0.07,0.07,0.07}{#1}}
\newcommand{\VerbatimStringTok}[1]{\textcolor[rgb]{0.13,0.47,0.30}{#1}}
\newcommand{\WarningTok}[1]{\textcolor[rgb]{0.37,0.37,0.37}{\textit{#1}}}

\providecommand{\tightlist}{%
  \setlength{\itemsep}{0pt}\setlength{\parskip}{0pt}}
\usepackage{longtable,booktabs,array}
\usepackage{calc} % for calculating minipage widths
% Correct order of tables after \paragraph or \subparagraph
\usepackage{etoolbox}
\makeatletter
\patchcmd\longtable{\par}{\if@noskipsec\mbox{}\fi\par}{}{}
\makeatother
% Allow footnotes in longtable head/foot
\IfFileExists{footnotehyper.sty}{\usepackage{footnotehyper}}{\usepackage{footnote}}
\makesavenoteenv{longtable}

\usepackage{graphicx}
\makeatletter
\newsavebox\pandoc@box
\newcommand*\pandocbounded[1]{% scales image to fit in text height/width
  \sbox\pandoc@box{#1}%
  \Gscale@div\@tempa{\textheight}{\dimexpr\ht\pandoc@box+\dp\pandoc@box\relax}%
  \Gscale@div\@tempb{\linewidth}{\wd\pandoc@box}%
  \ifdim\@tempb\p@<\@tempa\p@\let\@tempa\@tempb\fi% select the smaller of both
  \ifdim\@tempa\p@<\p@\scalebox{\@tempa}{\usebox\pandoc@box}%
  \else\usebox{\pandoc@box}%
  \fi%
}
% Set default figure placement to htbp
\def\fps@figure{htbp}
\makeatother







\usepackage{newtx}

\defaultfontfeatures{Scale=MatchLowercase}
\defaultfontfeatures[\rmfamily]{Ligatures=TeX,Scale=1}





\title{Métricas de performance en ML Python}


\shorttitle{PREDICCIÓN Y MÉTRICAS}


\usepackage{etoolbox}



\ccoppy{\textcopyright~2021}



\author{Edison Achalma}



\affiliation{
{Escuela Profesional de Economía, Universidad Nacional de San Cristóbal
de Huamanga}}




\leftheader{Achalma}

\date{2021-11-29}


\abstract{Este abstract será actualizado una vez que se complete el
contenido final del artículo. }

\keywords{keyword1, keyword2}

\authornote{\par{\addORCIDlink{Edison Achalma}{0000-0001-6996-3364}} 
\par{ }
\par{   El autor no tiene conflictos de interés que revelar.    Los
roles de autor se clasificaron utilizando la taxonomía de roles de
colaborador (CRediT; https://credit.niso.org/) de la siguiente
manera:  Edison Achalma:   conceptualización, metodología, análisis
formal, investigación, recursos, curación de
datos, redacción, visualización, supervisión, administración del
proyecto}
\par{La correspondencia relativa a este artículo debe dirigirse a Edison
Achalma, Escuela Profesional de Economía, Universidad Nacional de San
Cristóbal de Huamanga, Portal Independencia N°
57, Ayacucho, AYA 5001, Perú, Email: \href{mailto:elmer.achalma.09@unsch.edu.pe}{elmer.achalma.09@unsch.edu.pe}}
}

\makeatletter
\let\endoldlt\endlongtable
\def\endlongtable{
\hline
\endoldlt
}
\makeatother

\urlstyle{same}



\makeatletter
\@ifpackageloaded{caption}{}{\usepackage{caption}}
\AtBeginDocument{%
\ifdefined\contentsname
  \renewcommand*\contentsname{Tabla de contenidos}
\else
  \newcommand\contentsname{Tabla de contenidos}
\fi
\ifdefined\listfigurename
  \renewcommand*\listfigurename{Lista de Figuras}
\else
  \newcommand\listfigurename{Lista de Figuras}
\fi
\ifdefined\listtablename
  \renewcommand*\listtablename{Lista de Tablas}
\else
  \newcommand\listtablename{Lista de Tablas}
\fi
\ifdefined\figurename
  \renewcommand*\figurename{Figura}
\else
  \newcommand\figurename{Figura}
\fi
\ifdefined\tablename
  \renewcommand*\tablename{Tabla}
\else
  \newcommand\tablename{Tabla}
\fi
}
\@ifpackageloaded{float}{}{\usepackage{float}}
\floatstyle{ruled}
\@ifundefined{c@chapter}{\newfloat{codelisting}{h}{lop}}{\newfloat{codelisting}{h}{lop}[chapter]}
\floatname{codelisting}{Listado}
\newcommand*\listoflistings{\listof{codelisting}{Lista de Listados}}
\makeatother
\makeatletter
\makeatother
\makeatletter
\@ifpackageloaded{caption}{}{\usepackage{caption}}
\@ifpackageloaded{subcaption}{}{\usepackage{subcaption}}
\makeatother
\makeatletter
\@ifpackageloaded{fontawesome5}{}{\usepackage{fontawesome5}}
\makeatother

% From https://tex.stackexchange.com/a/645996/211326
%%% apa7 doesn't want to add appendix section titles in the toc
%%% let's make it do it
\makeatletter
\xpatchcmd{\appendix}
  {\par}
  {\addcontentsline{toc}{section}{\@currentlabelname}\par}
  {}{}
\makeatother

%% Disable longtable counter
%% https://tex.stackexchange.com/a/248395/211326

\usepackage{etoolbox}

\makeatletter
\patchcmd{\LT@caption}
  {\bgroup}
  {\bgroup\global\LTpatch@captiontrue}
  {}{}
\patchcmd{\longtable}
  {\par}
  {\par\global\LTpatch@captionfalse}
  {}{}
\apptocmd{\endlongtable}
  {\ifLTpatch@caption\else\addtocounter{table}{-1}\fi}
  {}{}
\newif\ifLTpatch@caption
\makeatother

\begin{document}

\maketitle


\hypertarget{toc}{}
\tableofcontents
\newpage
\section[Introduction]{Métricas de performance en ML Python}

\setcounter{secnumdepth}{3}

\setlength\LTleft{0pt}




En esta octava guía exploraremos el problema de predicción y las
métricas fundamentales para evaluar modelos de clasificación. Estos
conceptos son esenciales para evaluar modelos de machine learning
aplicados a problemas económicos como predicción de incumplimiento
crediticio, clasificación de empresas, predicción de pobreza, entre
otros.

\section{El problema de predicción}\label{el-problema-de-predicciuxf3n}

\subsection{Concepto fundamental}\label{concepto-fundamental}

En economía y ciencias sociales frecuentemente necesitamos predecir
eventos futuros basándonos en información histórica:

\begin{itemize}
\tightlist
\item
  ¿Un cliente pagará o no pagará su crédito?
\item
  ¿Una empresa quebrará o sobrevivirá?
\item
  ¿Un hogar saldrá o permanecerá en pobreza?
\item
  ¿Un estudiante completará o abandonará sus estudios?
\item
  ¿Un producto se venderá o no se venderá?
\end{itemize}

\begin{Shaded}
\begin{Highlighting}[]
\ImportTok{import}\NormalTok{ pandas }\ImportTok{as}\NormalTok{ pd}
\ImportTok{import}\NormalTok{ numpy }\ImportTok{as}\NormalTok{ np}
\ImportTok{import}\NormalTok{ matplotlib.pyplot }\ImportTok{as}\NormalTok{ plt}
\ImportTok{from}\NormalTok{ sklearn.linear\_model }\ImportTok{import}\NormalTok{ LogisticRegression}
\ImportTok{from}\NormalTok{ sklearn.metrics }\ImportTok{import}\NormalTok{ confusion\_matrix, classification\_report}
\ImportTok{from}\NormalTok{ sklearn.metrics }\ImportTok{import}\NormalTok{ roc\_curve, auc, precision\_recall\_curve}
\ImportTok{from}\NormalTok{ sklearn.model\_selection }\ImportTok{import}\NormalTok{ train\_test\_split}

\CommentTok{\# Configuración de gráficos}
\NormalTok{plt.style.use(}\StringTok{\textquotesingle{}seaborn{-}v0\_8{-}darkgrid\textquotesingle{}}\NormalTok{)}
\NormalTok{plt.rcParams[}\StringTok{\textquotesingle{}figure.figsize\textquotesingle{}}\NormalTok{] }\OperatorTok{=}\NormalTok{ (}\DecValTok{10}\NormalTok{, }\DecValTok{6}\NormalTok{)}
\end{Highlighting}
\end{Shaded}

\subsection{Flujo de decisión con
predicción}\label{flujo-de-decisiuxf3n-con-predicciuxf3n}

\begin{Shaded}
\begin{Highlighting}[]
\BuiltInTok{print}\NormalTok{(}\StringTok{"""}
\StringTok{FLUJO DE TOMA DE DECISIONES CON PREDICCIÓN}

\StringTok{1. RECOLECCIÓN DE DATOS}
\StringTok{   ├─ Variables explicativas (características)}
\StringTok{   └─ Variable objetivo (lo que queremos predecir)}
\StringTok{   }
\StringTok{2. ENTRENAMIENTO DEL MODELO}
\StringTok{   ├─ Aprender patrones de los datos históricos}
\StringTok{   └─ Establecer relación entre X (predictores) y Y (objetivo)}
\StringTok{   }
\StringTok{3. PREDICCIÓN}
\StringTok{   ├─ Aplicar modelo a nuevos casos}
\StringTok{   └─ Obtener probabilidades o clasificaciones}
\StringTok{   }
\StringTok{4. EVALUACIÓN}
\StringTok{   ├─ Medir qué tan buenas son las predicciones}
\StringTok{   └─ Usar métricas apropiadas}
\StringTok{   }
\StringTok{5. DECISIÓN}
\StringTok{   ├─ Usar predicción para tomar acción}
\StringTok{   └─ Ej: Aprobar/rechazar crédito, intervenir/no intervenir}
\StringTok{"""}\NormalTok{)}
\end{Highlighting}
\end{Shaded}

\section{Clasificación binaria}\label{clasificaciuxf3n-binaria}

\subsection{Problema de admisión
académica}\label{problema-de-admisiuxf3n-acaduxe9mica}

\begin{Shaded}
\begin{Highlighting}[]
\CommentTok{\# Crear dataset de ejemplo: admisión a programa de posgrado}
\NormalTok{np.random.seed(}\DecValTok{42}\NormalTok{)}

\CommentTok{\# Datos simulados}
\NormalTok{n\_estudiantes }\OperatorTok{=} \DecValTok{200}

\NormalTok{datos\_admision }\OperatorTok{=}\NormalTok{ pd.DataFrame(\{}
    \StringTok{\textquotesingle{}puntaje\_examen\textquotesingle{}}\NormalTok{: np.random.randint(}\DecValTok{500}\NormalTok{, }\DecValTok{800}\NormalTok{, n\_estudiantes),}
    \StringTok{\textquotesingle{}gpa\textquotesingle{}}\NormalTok{: np.random.uniform(}\FloatTok{2.0}\NormalTok{, }\FloatTok{4.0}\NormalTok{, n\_estudiantes),}
    \StringTok{\textquotesingle{}experiencia\_laboral\textquotesingle{}}\NormalTok{: np.random.randint(}\DecValTok{0}\NormalTok{, }\DecValTok{8}\NormalTok{, n\_estudiantes),}
    \StringTok{\textquotesingle{}publicaciones\textquotesingle{}}\NormalTok{: np.random.randint(}\DecValTok{0}\NormalTok{, }\DecValTok{5}\NormalTok{, n\_estudiantes)}
\NormalTok{\})}

\CommentTok{\# Generar variable de admisión (1 = admitido, 0 = no admitido)}
\CommentTok{\# Basado en una combinación de factores}
\NormalTok{score }\OperatorTok{=}\NormalTok{ (datos\_admision[}\StringTok{\textquotesingle{}puntaje\_examen\textquotesingle{}}\NormalTok{] }\OperatorTok{{-}} \DecValTok{500}\NormalTok{) }\OperatorTok{*} \FloatTok{0.002} \OperatorTok{+} \OperatorTok{\textbackslash{}}
\NormalTok{        datos\_admision[}\StringTok{\textquotesingle{}gpa\textquotesingle{}}\NormalTok{] }\OperatorTok{*} \FloatTok{0.15} \OperatorTok{+} \OperatorTok{\textbackslash{}}
\NormalTok{        datos\_admision[}\StringTok{\textquotesingle{}experiencia\_laboral\textquotesingle{}}\NormalTok{] }\OperatorTok{*} \FloatTok{0.08} \OperatorTok{+} \OperatorTok{\textbackslash{}}
\NormalTok{        datos\_admision[}\StringTok{\textquotesingle{}publicaciones\textquotesingle{}}\NormalTok{] }\OperatorTok{*} \FloatTok{0.10} \OperatorTok{+} \OperatorTok{\textbackslash{}}
\NormalTok{        np.random.normal(}\DecValTok{0}\NormalTok{, }\FloatTok{0.2}\NormalTok{, n\_estudiantes)}

\NormalTok{datos\_admision[}\StringTok{\textquotesingle{}admitido\textquotesingle{}}\NormalTok{] }\OperatorTok{=}\NormalTok{ (score }\OperatorTok{\textgreater{}} \FloatTok{1.0}\NormalTok{).astype(}\BuiltInTok{int}\NormalTok{)}

\BuiltInTok{print}\NormalTok{(}\StringTok{"Dataset de admisiones:"}\NormalTok{)}
\BuiltInTok{print}\NormalTok{(datos\_admision.head(}\DecValTok{10}\NormalTok{))}
\BuiltInTok{print}\NormalTok{(}\SpecialStringTok{f"}\CharTok{\textbackslash{}n}\SpecialStringTok{Distribución de admisiones:"}\NormalTok{)}
\BuiltInTok{print}\NormalTok{(datos\_admision[}\StringTok{\textquotesingle{}admitido\textquotesingle{}}\NormalTok{].value\_counts())}
\BuiltInTok{print}\NormalTok{(}\SpecialStringTok{f"}\CharTok{\textbackslash{}n}\SpecialStringTok{Tasa de admisión: }\SpecialCharTok{\{}\NormalTok{datos\_admision[}\StringTok{\textquotesingle{}admitido\textquotesingle{}}\NormalTok{]}\SpecialCharTok{.}\NormalTok{mean()}\OperatorTok{*}\DecValTok{100}\SpecialCharTok{:.1f\}}\SpecialStringTok{\%"}\NormalTok{)}
\end{Highlighting}
\end{Shaded}

\subsection{Entrenar modelo de
clasificación}\label{entrenar-modelo-de-clasificaciuxf3n}

\begin{Shaded}
\begin{Highlighting}[]
\CommentTok{\# Separar variables predictoras y objetivo}
\NormalTok{X }\OperatorTok{=}\NormalTok{ datos\_admision[[}\StringTok{\textquotesingle{}puntaje\_examen\textquotesingle{}}\NormalTok{, }\StringTok{\textquotesingle{}gpa\textquotesingle{}}\NormalTok{, }\StringTok{\textquotesingle{}experiencia\_laboral\textquotesingle{}}\NormalTok{, }\StringTok{\textquotesingle{}publicaciones\textquotesingle{}}\NormalTok{]]}
\NormalTok{y }\OperatorTok{=}\NormalTok{ datos\_admision[}\StringTok{\textquotesingle{}admitido\textquotesingle{}}\NormalTok{]}

\CommentTok{\# Dividir en conjunto de entrenamiento y prueba}
\NormalTok{X\_train, X\_test, y\_train, y\_test }\OperatorTok{=}\NormalTok{ train\_test\_split(}
\NormalTok{    X, y, test\_size}\OperatorTok{=}\FloatTok{0.3}\NormalTok{, random\_state}\OperatorTok{=}\DecValTok{42}
\NormalTok{)}

\BuiltInTok{print}\NormalTok{(}\SpecialStringTok{f"Tamaño conjunto entrenamiento: }\SpecialCharTok{\{}\BuiltInTok{len}\NormalTok{(X\_train)}\SpecialCharTok{\}}\SpecialStringTok{"}\NormalTok{)}
\BuiltInTok{print}\NormalTok{(}\SpecialStringTok{f"Tamaño conjunto prueba: }\SpecialCharTok{\{}\BuiltInTok{len}\NormalTok{(X\_test)}\SpecialCharTok{\}}\SpecialStringTok{"}\NormalTok{)}

\CommentTok{\# Entrenar modelo de regresión logística}
\NormalTok{modelo }\OperatorTok{=}\NormalTok{ LogisticRegression(random\_state}\OperatorTok{=}\DecValTok{42}\NormalTok{)}
\NormalTok{modelo.fit(X\_train, y\_train)}

\BuiltInTok{print}\NormalTok{(}\StringTok{"}\CharTok{\textbackslash{}n}\StringTok{Modelo entrenado exitosamente"}\NormalTok{)}
\BuiltInTok{print}\NormalTok{(}\SpecialStringTok{f"Coeficientes del modelo:"}\NormalTok{)}
\ControlFlowTok{for}\NormalTok{ variable, coef }\KeywordTok{in} \BuiltInTok{zip}\NormalTok{(X.columns, modelo.coef\_[}\DecValTok{0}\NormalTok{]):}
    \BuiltInTok{print}\NormalTok{(}\SpecialStringTok{f"  }\SpecialCharTok{\{}\NormalTok{variable}\SpecialCharTok{\}}\SpecialStringTok{: }\SpecialCharTok{\{}\NormalTok{coef}\SpecialCharTok{:.4f\}}\SpecialStringTok{"}\NormalTok{)}
\end{Highlighting}
\end{Shaded}

\section{Thresholding
(Umbralización)}\label{thresholding-umbralizaciuxf3n}

\subsection{Concepto de
probabilidades}\label{concepto-de-probabilidades}

\begin{Shaded}
\begin{Highlighting}[]
\CommentTok{\# Obtener probabilidades de predicción}
\NormalTok{probabilidades }\OperatorTok{=}\NormalTok{ modelo.predict\_proba(X\_test)}

\BuiltInTok{print}\NormalTok{(}\StringTok{"Probabilidades de predicción (primeras 10 observaciones):"}\NormalTok{)}
\BuiltInTok{print}\NormalTok{(}\StringTok{"}\CharTok{\textbackslash{}n}\StringTok{Índice | P(No admitido) | P(Admitido) | Predicción"}\NormalTok{)}
\BuiltInTok{print}\NormalTok{(}\StringTok{"{-}"} \OperatorTok{*} \DecValTok{55}\NormalTok{)}

\ControlFlowTok{for}\NormalTok{ i }\KeywordTok{in} \BuiltInTok{range}\NormalTok{(}\DecValTok{10}\NormalTok{):}
\NormalTok{    pred }\OperatorTok{=} \StringTok{"Admitido"} \ControlFlowTok{if}\NormalTok{ probabilidades[i, }\DecValTok{1}\NormalTok{] }\OperatorTok{\textgreater{}} \FloatTok{0.5} \ControlFlowTok{else} \StringTok{"No admitido"}
    \BuiltInTok{print}\NormalTok{(}\SpecialStringTok{f"}\SpecialCharTok{\{}\NormalTok{i}\SpecialCharTok{:6d\}}\SpecialStringTok{ | }\SpecialCharTok{\{}\NormalTok{probabilidades[i, }\DecValTok{0}\NormalTok{]}\SpecialCharTok{:14.4f\}}\SpecialStringTok{ | }\SpecialCharTok{\{}\NormalTok{probabilidades[i, }\DecValTok{1}\NormalTok{]}\SpecialCharTok{:11.4f\}}\SpecialStringTok{ | }\SpecialCharTok{\{}\NormalTok{pred}\SpecialCharTok{\}}\SpecialStringTok{"}\NormalTok{)}

\CommentTok{\# Crear DataFrame con resultados}
\NormalTok{resultados }\OperatorTok{=}\NormalTok{ pd.DataFrame(\{}
    \StringTok{\textquotesingle{}prob\_no\_admitido\textquotesingle{}}\NormalTok{: probabilidades[:, }\DecValTok{0}\NormalTok{],}
    \StringTok{\textquotesingle{}prob\_admitido\textquotesingle{}}\NormalTok{: probabilidades[:, }\DecValTok{1}\NormalTok{],}
    \StringTok{\textquotesingle{}real\textquotesingle{}}\NormalTok{: y\_test.values}
\NormalTok{\})}

\BuiltInTok{print}\NormalTok{(}\StringTok{"}\CharTok{\textbackslash{}n\textbackslash{}n}\StringTok{Estadísticas de probabilidades:"}\NormalTok{)}
\BuiltInTok{print}\NormalTok{(resultados[}\StringTok{\textquotesingle{}prob\_admitido\textquotesingle{}}\NormalTok{].describe())}
\end{Highlighting}
\end{Shaded}

\subsection{Efecto del umbral}\label{efecto-del-umbral}

\begin{Shaded}
\begin{Highlighting}[]
\CommentTok{\# Experimentar con diferentes umbrales}
\NormalTok{umbrales }\OperatorTok{=}\NormalTok{ [}\FloatTok{0.3}\NormalTok{, }\FloatTok{0.5}\NormalTok{, }\FloatTok{0.7}\NormalTok{, }\FloatTok{0.9}\NormalTok{]}

\BuiltInTok{print}\NormalTok{(}\StringTok{"EFECTO DEL UMBRAL EN LA CLASIFICACIÓN"}\NormalTok{)}
\BuiltInTok{print}\NormalTok{(}\StringTok{"="} \OperatorTok{*} \DecValTok{70}\NormalTok{)}

\ControlFlowTok{for}\NormalTok{ umbral }\KeywordTok{in}\NormalTok{ umbrales:}
    \CommentTok{\# Aplicar umbral}
\NormalTok{    predicciones }\OperatorTok{=}\NormalTok{ (probabilidades[:, }\DecValTok{1}\NormalTok{] }\OperatorTok{\textgreater{}}\NormalTok{ umbral).astype(}\BuiltInTok{int}\NormalTok{)}
    
    \CommentTok{\# Contar predicciones}
\NormalTok{    n\_positivos }\OperatorTok{=}\NormalTok{ predicciones.}\BuiltInTok{sum}\NormalTok{()}
\NormalTok{    tasa\_positivos }\OperatorTok{=}\NormalTok{ (n\_positivos }\OperatorTok{/} \BuiltInTok{len}\NormalTok{(predicciones)) }\OperatorTok{*} \DecValTok{100}
    
    \BuiltInTok{print}\NormalTok{(}\SpecialStringTok{f"}\CharTok{\textbackslash{}n}\SpecialStringTok{Umbral: }\SpecialCharTok{\{}\NormalTok{umbral}\SpecialCharTok{\}}\SpecialStringTok{"}\NormalTok{)}
    \BuiltInTok{print}\NormalTok{(}\SpecialStringTok{f"  Predicciones positivas: }\SpecialCharTok{\{}\NormalTok{n\_positivos}\SpecialCharTok{\}}\SpecialStringTok{ (}\SpecialCharTok{\{}\NormalTok{tasa\_positivos}\SpecialCharTok{:.1f\}}\SpecialStringTok{\%)"}\NormalTok{)}
    \BuiltInTok{print}\NormalTok{(}\SpecialStringTok{f"  Predicciones negativas: }\SpecialCharTok{\{}\BuiltInTok{len}\NormalTok{(predicciones) }\OperatorTok{{-}}\NormalTok{ n\_positivos}\SpecialCharTok{\}}\SpecialStringTok{"}\NormalTok{)}
\end{Highlighting}
\end{Shaded}

\subsection{Visualizar distribución de
probabilidades}\label{visualizar-distribuciuxf3n-de-probabilidades}

\begin{Shaded}
\begin{Highlighting}[]
\CommentTok{\# Crear gráfico de distribución de probabilidades}
\NormalTok{fig, axes }\OperatorTok{=}\NormalTok{ plt.subplots(}\DecValTok{1}\NormalTok{, }\DecValTok{2}\NormalTok{, figsize}\OperatorTok{=}\NormalTok{(}\DecValTok{14}\NormalTok{, }\DecValTok{5}\NormalTok{))}

\CommentTok{\# Histograma por clase real}
\ControlFlowTok{for}\NormalTok{ clase, nombre }\KeywordTok{in}\NormalTok{ [(}\DecValTok{0}\NormalTok{, }\StringTok{\textquotesingle{}No admitido\textquotesingle{}}\NormalTok{), (}\DecValTok{1}\NormalTok{, }\StringTok{\textquotesingle{}Admitido\textquotesingle{}}\NormalTok{)]:}
\NormalTok{    mask }\OperatorTok{=}\NormalTok{ resultados[}\StringTok{\textquotesingle{}real\textquotesingle{}}\NormalTok{] }\OperatorTok{==}\NormalTok{ clase}
\NormalTok{    axes[}\DecValTok{0}\NormalTok{].hist(resultados.loc[mask, }\StringTok{\textquotesingle{}prob\_admitido\textquotesingle{}}\NormalTok{], }
\NormalTok{                 bins}\OperatorTok{=}\DecValTok{20}\NormalTok{, alpha}\OperatorTok{=}\FloatTok{0.6}\NormalTok{, label}\OperatorTok{=}\NormalTok{nombre)}

\NormalTok{axes[}\DecValTok{0}\NormalTok{].axvline(x}\OperatorTok{=}\FloatTok{0.5}\NormalTok{, color}\OperatorTok{=}\StringTok{\textquotesingle{}red\textquotesingle{}}\NormalTok{, linestyle}\OperatorTok{=}\StringTok{\textquotesingle{}{-}{-}\textquotesingle{}}\NormalTok{, label}\OperatorTok{=}\StringTok{\textquotesingle{}Umbral = 0.5\textquotesingle{}}\NormalTok{)}
\NormalTok{axes[}\DecValTok{0}\NormalTok{].set\_xlabel(}\StringTok{\textquotesingle{}Probabilidad de admisión\textquotesingle{}}\NormalTok{)}
\NormalTok{axes[}\DecValTok{0}\NormalTok{].set\_ylabel(}\StringTok{\textquotesingle{}Frecuencia\textquotesingle{}}\NormalTok{)}
\NormalTok{axes[}\DecValTok{0}\NormalTok{].set\_title(}\StringTok{\textquotesingle{}Distribución de probabilidades por clase real\textquotesingle{}}\NormalTok{)}
\NormalTok{axes[}\DecValTok{0}\NormalTok{].legend()}
\NormalTok{axes[}\DecValTok{0}\NormalTok{].grid(}\VariableTok{True}\NormalTok{, alpha}\OperatorTok{=}\FloatTok{0.3}\NormalTok{)}

\CommentTok{\# Scatter plot}
\NormalTok{colores }\OperatorTok{=}\NormalTok{ [}\StringTok{\textquotesingle{}red\textquotesingle{}} \ControlFlowTok{if}\NormalTok{ x }\OperatorTok{==} \DecValTok{0} \ControlFlowTok{else} \StringTok{\textquotesingle{}green\textquotesingle{}} \ControlFlowTok{for}\NormalTok{ x }\KeywordTok{in}\NormalTok{ resultados[}\StringTok{\textquotesingle{}real\textquotesingle{}}\NormalTok{]]}
\NormalTok{axes[}\DecValTok{1}\NormalTok{].scatter(}\BuiltInTok{range}\NormalTok{(}\BuiltInTok{len}\NormalTok{(resultados)), }
\NormalTok{                resultados[}\StringTok{\textquotesingle{}prob\_admitido\textquotesingle{}}\NormalTok{],}
\NormalTok{                c}\OperatorTok{=}\NormalTok{colores, alpha}\OperatorTok{=}\FloatTok{0.6}\NormalTok{)}
\NormalTok{axes[}\DecValTok{1}\NormalTok{].axhline(y}\OperatorTok{=}\FloatTok{0.5}\NormalTok{, color}\OperatorTok{=}\StringTok{\textquotesingle{}blue\textquotesingle{}}\NormalTok{, linestyle}\OperatorTok{=}\StringTok{\textquotesingle{}{-}{-}\textquotesingle{}}\NormalTok{, label}\OperatorTok{=}\StringTok{\textquotesingle{}Umbral = 0.5\textquotesingle{}}\NormalTok{)}
\NormalTok{axes[}\DecValTok{1}\NormalTok{].set\_xlabel(}\StringTok{\textquotesingle{}Índice de observación\textquotesingle{}}\NormalTok{)}
\NormalTok{axes[}\DecValTok{1}\NormalTok{].set\_ylabel(}\StringTok{\textquotesingle{}Probabilidad de admisión\textquotesingle{}}\NormalTok{)}
\NormalTok{axes[}\DecValTok{1}\NormalTok{].set\_title(}\StringTok{\textquotesingle{}Probabilidades predichas vs clase real\textquotesingle{}}\NormalTok{)}
\NormalTok{axes[}\DecValTok{1}\NormalTok{].legend([}\StringTok{\textquotesingle{}Umbral 0.5\textquotesingle{}}\NormalTok{, }\StringTok{\textquotesingle{}No admitido\textquotesingle{}}\NormalTok{, }\StringTok{\textquotesingle{}Admitido\textquotesingle{}}\NormalTok{])}
\NormalTok{axes[}\DecValTok{1}\NormalTok{].grid(}\VariableTok{True}\NormalTok{, alpha}\OperatorTok{=}\FloatTok{0.3}\NormalTok{)}

\NormalTok{plt.tight\_layout()}
\NormalTok{plt.savefig(}\StringTok{\textquotesingle{}distribucion\_probabilidades.png\textquotesingle{}}\NormalTok{, dpi}\OperatorTok{=}\DecValTok{300}\NormalTok{, bbox\_inches}\OperatorTok{=}\StringTok{\textquotesingle{}tight\textquotesingle{}}\NormalTok{)}
\BuiltInTok{print}\NormalTok{(}\StringTok{"}\CharTok{\textbackslash{}n}\StringTok{Gráfico guardado como \textquotesingle{}distribucion\_probabilidades.png\textquotesingle{}"}\NormalTok{)}
\end{Highlighting}
\end{Shaded}

\section{Matriz de confusión}\label{matriz-de-confusiuxf3n}

\subsection{Concepto y estructura}\label{concepto-y-estructura}

La matriz de confusión es una tabla que describe el desempeño de un
modelo de clasificación.

\begin{Shaded}
\begin{Highlighting}[]
\BuiltInTok{print}\NormalTok{(}\StringTok{"""}
\StringTok{MATRIZ DE CONFUSIÓN {-} ESTRUCTURA}

\StringTok{                    Predicción}
\StringTok{                 Negativo  Positivo}
\StringTok{               ┌──────────┬──────────┐}
\StringTok{Real  Negativo │    TN    │    FP    │}
\StringTok{               ├──────────┼──────────┤}
\StringTok{      Positivo │    FN    │    TP    │}
\StringTok{               └──────────┴──────────┘}

\StringTok{Donde:}
\StringTok{{-} TN (True Negative): Negativos correctamente identificados}
\StringTok{{-} FP (False Positive): Negativos incorrectamente clasificados como positivos}
\StringTok{{-} FN (False Negative): Positivos incorrectamente clasificados como negativos}
\StringTok{{-} TP (True Positive): Positivos correctamente identificados}
\StringTok{"""}\NormalTok{)}
\end{Highlighting}
\end{Shaded}

\subsection{Calcular matriz de
confusión}\label{calcular-matriz-de-confusiuxf3n}

\begin{Shaded}
\begin{Highlighting}[]
\CommentTok{\# Predicciones con umbral 0.5}
\NormalTok{y\_pred }\OperatorTok{=}\NormalTok{ modelo.predict(X\_test)}

\CommentTok{\# Calcular matriz de confusión}
\NormalTok{cm }\OperatorTok{=}\NormalTok{ confusion\_matrix(y\_test, y\_pred)}

\BuiltInTok{print}\NormalTok{(}\StringTok{"MATRIZ DE CONFUSIÓN"}\NormalTok{)}
\BuiltInTok{print}\NormalTok{(}\StringTok{"="} \OperatorTok{*} \DecValTok{50}\NormalTok{)}
\BuiltInTok{print}\NormalTok{(}\SpecialStringTok{f"}\CharTok{\textbackslash{}n}\SpecialStringTok{                 Predicción"}\NormalTok{)}
\BuiltInTok{print}\NormalTok{(}\SpecialStringTok{f"              No admitido  Admitido"}\NormalTok{)}
\BuiltInTok{print}\NormalTok{(}\SpecialStringTok{f"Real No admitido    }\SpecialCharTok{\{}\NormalTok{cm[}\DecValTok{0}\NormalTok{,}\DecValTok{0}\NormalTok{]}\SpecialCharTok{:4d\}}\SpecialStringTok{       }\SpecialCharTok{\{}\NormalTok{cm[}\DecValTok{0}\NormalTok{,}\DecValTok{1}\NormalTok{]}\SpecialCharTok{:4d\}}\SpecialStringTok{"}\NormalTok{)}
\BuiltInTok{print}\NormalTok{(}\SpecialStringTok{f"     Admitido       }\SpecialCharTok{\{}\NormalTok{cm[}\DecValTok{1}\NormalTok{,}\DecValTok{0}\NormalTok{]}\SpecialCharTok{:4d\}}\SpecialStringTok{       }\SpecialCharTok{\{}\NormalTok{cm[}\DecValTok{1}\NormalTok{,}\DecValTok{1}\NormalTok{]}\SpecialCharTok{:4d\}}\SpecialStringTok{"}\NormalTok{)}

\CommentTok{\# Extraer componentes}
\NormalTok{TN, FP, FN, TP }\OperatorTok{=}\NormalTok{ cm.ravel()}

\BuiltInTok{print}\NormalTok{(}\SpecialStringTok{f"}\CharTok{\textbackslash{}n\textbackslash{}n}\SpecialStringTok{Componentes de la matriz:"}\NormalTok{)}
\BuiltInTok{print}\NormalTok{(}\SpecialStringTok{f"  Verdaderos Negativos (TN): }\SpecialCharTok{\{}\NormalTok{TN}\SpecialCharTok{\}}\SpecialStringTok{"}\NormalTok{)}
\BuiltInTok{print}\NormalTok{(}\SpecialStringTok{f"  Falsos Positivos (FP): }\SpecialCharTok{\{}\NormalTok{FP}\SpecialCharTok{\}}\SpecialStringTok{"}\NormalTok{)}
\BuiltInTok{print}\NormalTok{(}\SpecialStringTok{f"  Falsos Negativos (FN): }\SpecialCharTok{\{}\NormalTok{FN}\SpecialCharTok{\}}\SpecialStringTok{"}\NormalTok{)}
\BuiltInTok{print}\NormalTok{(}\SpecialStringTok{f"  Verdaderos Positivos (TP): }\SpecialCharTok{\{}\NormalTok{TP}\SpecialCharTok{\}}\SpecialStringTok{"}\NormalTok{)}
\end{Highlighting}
\end{Shaded}

\subsection{Función personalizada para matriz de
confusión}\label{funciuxf3n-personalizada-para-matriz-de-confusiuxf3n}

\begin{Shaded}
\begin{Highlighting}[]
\KeywordTok{def}\NormalTok{ calcular\_matriz\_confusion(y\_real, y\_pred):}
    \CommentTok{"""}
\CommentTok{    Calcula matriz de confusión y sus componentes}
\CommentTok{    """}
    \CommentTok{\# Convertir a arrays booleanos}
\NormalTok{    y\_real }\OperatorTok{=}\NormalTok{ np.array(y\_real).astype(}\BuiltInTok{bool}\NormalTok{)}
\NormalTok{    y\_pred }\OperatorTok{=}\NormalTok{ np.array(y\_pred).astype(}\BuiltInTok{bool}\NormalTok{)}
    
    \CommentTok{\# Calcular componentes}
\NormalTok{    TP }\OperatorTok{=}\NormalTok{ (y\_real }\OperatorTok{\&}\NormalTok{ y\_pred).}\BuiltInTok{sum}\NormalTok{()}
\NormalTok{    TN }\OperatorTok{=}\NormalTok{ (}\OperatorTok{\textasciitilde{}}\NormalTok{y\_real }\OperatorTok{\&} \OperatorTok{\textasciitilde{}}\NormalTok{y\_pred).}\BuiltInTok{sum}\NormalTok{()}
\NormalTok{    FP }\OperatorTok{=}\NormalTok{ (}\OperatorTok{\textasciitilde{}}\NormalTok{y\_real }\OperatorTok{\&}\NormalTok{ y\_pred).}\BuiltInTok{sum}\NormalTok{()}
\NormalTok{    FN }\OperatorTok{=}\NormalTok{ (y\_real }\OperatorTok{\&} \OperatorTok{\textasciitilde{}}\NormalTok{y\_pred).}\BuiltInTok{sum}\NormalTok{()}
    
    \CommentTok{\# Crear matriz}
\NormalTok{    matriz }\OperatorTok{=}\NormalTok{ np.array([[TN, FP],}
\NormalTok{                       [FN, TP]])}
    
    \ControlFlowTok{return}\NormalTok{ matriz, TP, TN, FP, FN}

\CommentTok{\# Usar función personalizada}
\NormalTok{matriz, TP, TN, FP, FN }\OperatorTok{=}\NormalTok{ calcular\_matriz\_confusion(y\_test, y\_pred)}

\BuiltInTok{print}\NormalTok{(}\StringTok{"Matriz de confusión (función personalizada):"}\NormalTok{)}
\BuiltInTok{print}\NormalTok{(matriz)}
\BuiltInTok{print}\NormalTok{(}\SpecialStringTok{f"}\CharTok{\textbackslash{}n}\SpecialStringTok{Verificación: TP=}\SpecialCharTok{\{}\NormalTok{TP}\SpecialCharTok{\}}\SpecialStringTok{, TN=}\SpecialCharTok{\{}\NormalTok{TN}\SpecialCharTok{\}}\SpecialStringTok{, FP=}\SpecialCharTok{\{}\NormalTok{FP}\SpecialCharTok{\}}\SpecialStringTok{, FN=}\SpecialCharTok{\{}\NormalTok{FN}\SpecialCharTok{\}}\SpecialStringTok{"}\NormalTok{)}
\end{Highlighting}
\end{Shaded}

\subsection{Visualizar matriz de
confusión}\label{visualizar-matriz-de-confusiuxf3n}

\begin{Shaded}
\begin{Highlighting}[]
\KeywordTok{def}\NormalTok{ graficar\_matriz\_confusion(cm, nombres\_clases}\OperatorTok{=}\NormalTok{[}\StringTok{\textquotesingle{}No\textquotesingle{}}\NormalTok{, }\StringTok{\textquotesingle{}Sí\textquotesingle{}}\NormalTok{]):}
    \CommentTok{"""}
\CommentTok{    Visualiza la matriz de confusión}
\CommentTok{    """}
\NormalTok{    fig, ax }\OperatorTok{=}\NormalTok{ plt.subplots(figsize}\OperatorTok{=}\NormalTok{(}\DecValTok{8}\NormalTok{, }\DecValTok{6}\NormalTok{))}
    
\NormalTok{    im }\OperatorTok{=}\NormalTok{ ax.imshow(cm, interpolation}\OperatorTok{=}\StringTok{\textquotesingle{}nearest\textquotesingle{}}\NormalTok{, cmap}\OperatorTok{=}\NormalTok{plt.cm.Blues)}
\NormalTok{    ax.figure.colorbar(im, ax}\OperatorTok{=}\NormalTok{ax)}
    
\NormalTok{    ax.}\BuiltInTok{set}\NormalTok{(xticks}\OperatorTok{=}\NormalTok{np.arange(cm.shape[}\DecValTok{1}\NormalTok{]),}
\NormalTok{           yticks}\OperatorTok{=}\NormalTok{np.arange(cm.shape[}\DecValTok{0}\NormalTok{]),}
\NormalTok{           xticklabels}\OperatorTok{=}\NormalTok{nombres\_clases,}
\NormalTok{           yticklabels}\OperatorTok{=}\NormalTok{nombres\_clases,}
\NormalTok{           ylabel}\OperatorTok{=}\StringTok{\textquotesingle{}Clase Real\textquotesingle{}}\NormalTok{,}
\NormalTok{           xlabel}\OperatorTok{=}\StringTok{\textquotesingle{}Clase Predicha\textquotesingle{}}\NormalTok{,}
\NormalTok{           title}\OperatorTok{=}\StringTok{\textquotesingle{}Matriz de Confusión\textquotesingle{}}\NormalTok{)}
    
    \CommentTok{\# Anotar celdas con valores}
\NormalTok{    fmt }\OperatorTok{=} \StringTok{\textquotesingle{}d\textquotesingle{}}
\NormalTok{    thresh }\OperatorTok{=}\NormalTok{ cm.}\BuiltInTok{max}\NormalTok{() }\OperatorTok{/} \FloatTok{2.}
    \ControlFlowTok{for}\NormalTok{ i }\KeywordTok{in} \BuiltInTok{range}\NormalTok{(cm.shape[}\DecValTok{0}\NormalTok{]):}
        \ControlFlowTok{for}\NormalTok{ j }\KeywordTok{in} \BuiltInTok{range}\NormalTok{(cm.shape[}\DecValTok{1}\NormalTok{]):}
\NormalTok{            ax.text(j, i, }\BuiltInTok{format}\NormalTok{(cm[i, j], fmt),}
\NormalTok{                   ha}\OperatorTok{=}\StringTok{"center"}\NormalTok{, va}\OperatorTok{=}\StringTok{"center"}\NormalTok{,}
\NormalTok{                   color}\OperatorTok{=}\StringTok{"white"} \ControlFlowTok{if}\NormalTok{ cm[i, j] }\OperatorTok{\textgreater{}}\NormalTok{ thresh }\ControlFlowTok{else} \StringTok{"black"}\NormalTok{,}
\NormalTok{                   fontsize}\OperatorTok{=}\DecValTok{20}\NormalTok{)}
    
\NormalTok{    plt.tight\_layout()}
    \ControlFlowTok{return}\NormalTok{ fig}

\CommentTok{\# Graficar}
\NormalTok{fig }\OperatorTok{=}\NormalTok{ graficar\_matriz\_confusion(cm, nombres\_clases}\OperatorTok{=}\NormalTok{[}\StringTok{\textquotesingle{}No admitido\textquotesingle{}}\NormalTok{, }\StringTok{\textquotesingle{}Admitido\textquotesingle{}}\NormalTok{])}
\NormalTok{plt.savefig(}\StringTok{\textquotesingle{}matriz\_confusion.png\textquotesingle{}}\NormalTok{, dpi}\OperatorTok{=}\DecValTok{300}\NormalTok{, bbox\_inches}\OperatorTok{=}\StringTok{\textquotesingle{}tight\textquotesingle{}}\NormalTok{)}
\BuiltInTok{print}\NormalTok{(}\StringTok{"Matriz de confusión guardada como \textquotesingle{}matriz\_confusion.png\textquotesingle{}"}\NormalTok{)}
\end{Highlighting}
\end{Shaded}

\section{Métricas de
clasificación}\label{muxe9tricas-de-clasificaciuxf3n}

\subsection{Resumen de métricas
principales}\label{resumen-de-muxe9tricas-principales}

\begin{Shaded}
\begin{Highlighting}[]
\BuiltInTok{print}\NormalTok{(}\StringTok{"""}
\StringTok{MÉTRICAS DE CLASIFICACIÓN}

\StringTok{1. ACCURACY (Exactitud)}
\StringTok{   {-} Proporción total de predicciones correctas}
\StringTok{   {-} Fórmula: (TP + TN) / (TP + TN + FP + FN)}
\StringTok{   {-} Cuándo usar: Clases balanceadas}
\StringTok{   }
\StringTok{2. PRECISION (Precisión)}
\StringTok{   {-} De las predicciones positivas, cuántas son correctas}
\StringTok{   {-} Fórmula: TP / (TP + FP)}
\StringTok{   {-} Cuándo usar: Costo alto de falsos positivos}
\StringTok{   {-} Ejemplo: Diagnóstico médico (evitar alarmas falsas)}
\StringTok{   }
\StringTok{3. RECALL (Sensibilidad/Exhaustividad)}
\StringTok{   {-} De los casos positivos reales, cuántos detectamos}
\StringTok{   {-} Fórmula: TP / (TP + FN)}
\StringTok{   {-} Cuándo usar: Costo alto de falsos negativos}
\StringTok{   {-} Ejemplo: Detección de fraude (no perder casos reales)}
\StringTok{   }
\StringTok{4. F1{-}SCORE}
\StringTok{   {-} Media armónica de precision y recall}
\StringTok{   {-} Fórmula: 2 * (Precision * Recall) / (Precision + Recall)}
\StringTok{   {-} Cuándo usar: Balance entre precision y recall}
\StringTok{   }
\StringTok{5. SPECIFICITY (Especificidad)}
\StringTok{   {-} De los casos negativos reales, cuántos detectamos}
\StringTok{   {-} Fórmula: TN / (TN + FP)}
\StringTok{"""}\NormalTok{)}
\end{Highlighting}
\end{Shaded}

\section{Accuracy (Exactitud)}\label{accuracy-exactitud}

\subsection{Cálculo e
interpretación}\label{cuxe1lculo-e-interpretaciuxf3n}

\begin{Shaded}
\begin{Highlighting}[]
\KeywordTok{def}\NormalTok{ calcular\_accuracy(y\_real, y\_pred):}
    \CommentTok{"""Calcula accuracy"""}
\NormalTok{    \_, TP, TN, FP, FN }\OperatorTok{=}\NormalTok{ calcular\_matriz\_confusion(y\_real, y\_pred)}
\NormalTok{    accuracy }\OperatorTok{=}\NormalTok{ (TP }\OperatorTok{+}\NormalTok{ TN) }\OperatorTok{/}\NormalTok{ (TP }\OperatorTok{+}\NormalTok{ TN }\OperatorTok{+}\NormalTok{ FP }\OperatorTok{+}\NormalTok{ FN)}
    \ControlFlowTok{return}\NormalTok{ accuracy}

\CommentTok{\# Calcular accuracy}
\NormalTok{accuracy }\OperatorTok{=}\NormalTok{ calcular\_accuracy(y\_test, y\_pred)}
\NormalTok{accuracy\_sklearn }\OperatorTok{=}\NormalTok{ modelo.score(X\_test, y\_test)}

\BuiltInTok{print}\NormalTok{(}\SpecialStringTok{f"Accuracy (función propia): }\SpecialCharTok{\{}\NormalTok{accuracy}\SpecialCharTok{:.4f\}}\SpecialStringTok{ (}\SpecialCharTok{\{}\NormalTok{accuracy}\OperatorTok{*}\DecValTok{100}\SpecialCharTok{:.2f\}}\SpecialStringTok{\%)"}\NormalTok{)}
\BuiltInTok{print}\NormalTok{(}\SpecialStringTok{f"Accuracy (sklearn): }\SpecialCharTok{\{}\NormalTok{accuracy\_sklearn}\SpecialCharTok{:.4f\}}\SpecialStringTok{ (}\SpecialCharTok{\{}\NormalTok{accuracy\_sklearn}\OperatorTok{*}\DecValTok{100}\SpecialCharTok{:.2f\}}\SpecialStringTok{\%)"}\NormalTok{)}

\BuiltInTok{print}\NormalTok{(}\SpecialStringTok{f"}\CharTok{\textbackslash{}n}\SpecialStringTok{Interpretación:"}\NormalTok{)}
\BuiltInTok{print}\NormalTok{(}\SpecialStringTok{f"El modelo clasifica correctamente el }\SpecialCharTok{\{}\NormalTok{accuracy}\OperatorTok{*}\DecValTok{100}\SpecialCharTok{:.1f\}}\SpecialStringTok{\% de los casos"}\NormalTok{)}
\end{Highlighting}
\end{Shaded}

\subsection{Limitaciones del accuracy}\label{limitaciones-del-accuracy}

\begin{Shaded}
\begin{Highlighting}[]
\CommentTok{\# Ejemplo con dataset desbalanceado}
\BuiltInTok{print}\NormalTok{(}\StringTok{"}\CharTok{\textbackslash{}n}\StringTok{EJEMPLO: LIMITACIÓN DEL ACCURACY CON CLASES DESBALANCEADAS"}\NormalTok{)}
\BuiltInTok{print}\NormalTok{(}\StringTok{"="} \OperatorTok{*} \DecValTok{70}\NormalTok{)}

\CommentTok{\# Dataset muy desbalanceado (95\% negativos, 5\% positivos)}
\NormalTok{n }\OperatorTok{=} \DecValTok{1000}
\NormalTok{y\_desbalanceado }\OperatorTok{=}\NormalTok{ np.array([}\DecValTok{0}\NormalTok{]}\OperatorTok{*}\DecValTok{950} \OperatorTok{+}\NormalTok{ [}\DecValTok{1}\NormalTok{]}\OperatorTok{*}\DecValTok{50}\NormalTok{)}

\CommentTok{\# Modelo que siempre predice negativo}
\NormalTok{y\_pred\_siempre\_negativo }\OperatorTok{=}\NormalTok{ np.zeros(n)}

\NormalTok{acc }\OperatorTok{=}\NormalTok{ calcular\_accuracy(y\_desbalanceado, y\_pred\_siempre\_negativo)}
\BuiltInTok{print}\NormalTok{(}\SpecialStringTok{f"}\CharTok{\textbackslash{}n}\SpecialStringTok{Modelo que predice siempre \textquotesingle{}No\textquotesingle{}:"}\NormalTok{)}
\BuiltInTok{print}\NormalTok{(}\SpecialStringTok{f"Accuracy: }\SpecialCharTok{\{}\NormalTok{acc}\SpecialCharTok{:.4f\}}\SpecialStringTok{ (}\SpecialCharTok{\{}\NormalTok{acc}\OperatorTok{*}\DecValTok{100}\SpecialCharTok{:.1f\}}\SpecialStringTok{\%)"}\NormalTok{)}
\BuiltInTok{print}\NormalTok{(}\StringTok{"}\CharTok{\textbackslash{}n}\StringTok{¡Accuracy alto pero modelo inútil!"}\NormalTok{)}
\BuiltInTok{print}\NormalTok{(}\StringTok{"No detecta ningún caso positivo."}\NormalTok{)}

\CommentTok{\# Matriz de confusión}
\NormalTok{cm\_desbalanceado }\OperatorTok{=}\NormalTok{ confusion\_matrix(y\_desbalanceado, y\_pred\_siempre\_negativo)}
\BuiltInTok{print}\NormalTok{(}\SpecialStringTok{f"}\CharTok{\textbackslash{}n}\SpecialStringTok{Matriz de confusión:"}\NormalTok{)}
\BuiltInTok{print}\NormalTok{(cm\_desbalanceado)}
\BuiltInTok{print}\NormalTok{(}\StringTok{"}\CharTok{\textbackslash{}n}\StringTok{Observación: TP = 0 (no detecta ningún positivo)"}\NormalTok{)}
\end{Highlighting}
\end{Shaded}

\section{Precision (Precisión)}\label{precision-precisiuxf3n}

\subsection{Cálculo e
interpretación}\label{cuxe1lculo-e-interpretaciuxf3n-1}

\begin{Shaded}
\begin{Highlighting}[]
\KeywordTok{def}\NormalTok{ calcular\_precision(y\_real, y\_pred):}
    \CommentTok{"""Calcula precision"""}
\NormalTok{    \_, TP, TN, FP, FN }\OperatorTok{=}\NormalTok{ calcular\_matriz\_confusion(y\_real, y\_pred)}
    
    \ControlFlowTok{if}\NormalTok{ TP }\OperatorTok{+}\NormalTok{ FP }\OperatorTok{==} \DecValTok{0}\NormalTok{:}
        \ControlFlowTok{return} \FloatTok{0.0}
    
\NormalTok{    precision }\OperatorTok{=}\NormalTok{ TP }\OperatorTok{/}\NormalTok{ (TP }\OperatorTok{+}\NormalTok{ FP)}
    \ControlFlowTok{return}\NormalTok{ precision}

\CommentTok{\# Calcular precision}
\NormalTok{precision }\OperatorTok{=}\NormalTok{ calcular\_precision(y\_test, y\_pred)}

\BuiltInTok{print}\NormalTok{(}\SpecialStringTok{f"Precision: }\SpecialCharTok{\{}\NormalTok{precision}\SpecialCharTok{:.4f\}}\SpecialStringTok{ (}\SpecialCharTok{\{}\NormalTok{precision}\OperatorTok{*}\DecValTok{100}\SpecialCharTok{:.2f\}}\SpecialStringTok{\%)"}\NormalTok{)}
\BuiltInTok{print}\NormalTok{(}\SpecialStringTok{f"}\CharTok{\textbackslash{}n}\SpecialStringTok{Interpretación:"}\NormalTok{)}
\BuiltInTok{print}\NormalTok{(}\SpecialStringTok{f"De los estudiantes que el modelo predice como \textquotesingle{}admitidos\textquotesingle{},"}\NormalTok{)}
\BuiltInTok{print}\NormalTok{(}\SpecialStringTok{f"el }\SpecialCharTok{\{}\NormalTok{precision}\OperatorTok{*}\DecValTok{100}\SpecialCharTok{:.1f\}}\SpecialStringTok{\% realmente son admitidos."}\NormalTok{)}
\BuiltInTok{print}\NormalTok{(}\SpecialStringTok{f"}\CharTok{\textbackslash{}n}\SpecialStringTok{Falsos positivos: }\SpecialCharTok{\{}\NormalTok{FP}\SpecialCharTok{\}}\SpecialStringTok{ estudiantes"}\NormalTok{)}
\BuiltInTok{print}\NormalTok{(}\StringTok{"(Predichos como admitidos pero en realidad no lo fueron)"}\NormalTok{)}
\end{Highlighting}
\end{Shaded}

\subsection{Ejemplo económico: Aprobación de
créditos}\label{ejemplo-econuxf3mico-aprobaciuxf3n-de-cruxe9ditos}

\begin{Shaded}
\begin{Highlighting}[]
\BuiltInTok{print}\NormalTok{(}\StringTok{"}\CharTok{\textbackslash{}n\textbackslash{}n}\StringTok{EJEMPLO ECONÓMICO: APROBACIÓN DE CRÉDITOS"}\NormalTok{)}
\BuiltInTok{print}\NormalTok{(}\StringTok{"="} \OperatorTok{*} \DecValTok{70}\NormalTok{)}

\CommentTok{\# Simular decisiones de crédito}
\NormalTok{n\_creditos }\OperatorTok{=} \DecValTok{1000}
\NormalTok{y\_real\_pagara }\OperatorTok{=}\NormalTok{ np.random.choice([}\DecValTok{0}\NormalTok{, }\DecValTok{1}\NormalTok{], n\_creditos, p}\OperatorTok{=}\NormalTok{[}\FloatTok{0.2}\NormalTok{, }\FloatTok{0.8}\NormalTok{])  }\CommentTok{\# 80\% pagan}
\NormalTok{y\_pred\_aprobar }\OperatorTok{=}\NormalTok{ np.random.choice([}\DecValTok{0}\NormalTok{, }\DecValTok{1}\NormalTok{], n\_creditos, p}\OperatorTok{=}\NormalTok{[}\FloatTok{0.15}\NormalTok{, }\FloatTok{0.85}\NormalTok{])  }\CommentTok{\# Modelo aprueba 85\%}

\CommentTok{\# Calcular métricas}
\NormalTok{\_, TP, TN, FP, FN }\OperatorTok{=}\NormalTok{ calcular\_matriz\_confusion(y\_real\_pagara, y\_pred\_aprobar)}
\NormalTok{precision\_credito }\OperatorTok{=}\NormalTok{ calcular\_precision(y\_real\_pagara, y\_pred\_aprobar)}

\BuiltInTok{print}\NormalTok{(}\SpecialStringTok{f"}\CharTok{\textbackslash{}n}\SpecialStringTok{Resultados de }\SpecialCharTok{\{}\NormalTok{n\_creditos}\SpecialCharTok{\}}\SpecialStringTok{ solicitudes:"}\NormalTok{)}
\BuiltInTok{print}\NormalTok{(}\SpecialStringTok{f"  Créditos aprobados por modelo: }\SpecialCharTok{\{}\NormalTok{TP }\OperatorTok{+}\NormalTok{ FP}\SpecialCharTok{\}}\SpecialStringTok{"}\NormalTok{)}
\BuiltInTok{print}\NormalTok{(}\SpecialStringTok{f"  De estos, pagarán: }\SpecialCharTok{\{}\NormalTok{TP}\SpecialCharTok{\}}\SpecialStringTok{ (}\SpecialCharTok{\{}\NormalTok{precision\_credito}\OperatorTok{*}\DecValTok{100}\SpecialCharTok{:.1f\}}\SpecialStringTok{\%)"}\NormalTok{)}
\BuiltInTok{print}\NormalTok{(}\SpecialStringTok{f"  De estos, NO pagarán: }\SpecialCharTok{\{}\NormalTok{FP}\SpecialCharTok{\}}\SpecialStringTok{ (Pérdida por falsos positivos)"}\NormalTok{)}

\CommentTok{\# Calcular pérdida}
\NormalTok{monto\_promedio }\OperatorTok{=} \DecValTok{10000}
\NormalTok{perdida\_por\_fp }\OperatorTok{=}\NormalTok{ FP }\OperatorTok{*}\NormalTok{ monto\_promedio}

\BuiltInTok{print}\NormalTok{(}\SpecialStringTok{f"}\CharTok{\textbackslash{}n\textbackslash{}n}\SpecialStringTok{Análisis financiero:"}\NormalTok{)}
\BuiltInTok{print}\NormalTok{(}\SpecialStringTok{f"  Monto promedio por crédito: S/ }\SpecialCharTok{\{}\NormalTok{monto\_promedio}\SpecialCharTok{:,\}}\SpecialStringTok{"}\NormalTok{)}
\BuiltInTok{print}\NormalTok{(}\SpecialStringTok{f"  Falsos positivos: }\SpecialCharTok{\{}\NormalTok{FP}\SpecialCharTok{\}}\SpecialStringTok{"}\NormalTok{)}
\BuiltInTok{print}\NormalTok{(}\SpecialStringTok{f"  Pérdida estimada: S/ }\SpecialCharTok{\{}\NormalTok{perdida\_por\_fp}\SpecialCharTok{:,\}}\SpecialStringTok{"}\NormalTok{)}
\BuiltInTok{print}\NormalTok{(}\SpecialStringTok{f"}\CharTok{\textbackslash{}n}\SpecialStringTok{Conlusión: Alta precision reduce pérdidas por morosidad"}\NormalTok{)}
\end{Highlighting}
\end{Shaded}

\section{Recall (Sensibilidad)}\label{recall-sensibilidad}

\subsection{Cálculo e
interpretación}\label{cuxe1lculo-e-interpretaciuxf3n-2}

\begin{Shaded}
\begin{Highlighting}[]
\KeywordTok{def}\NormalTok{ calcular\_recall(y\_real, y\_pred):}
    \CommentTok{"""Calcula recall (sensibilidad)"""}
\NormalTok{    \_, TP, TN, FP, FN }\OperatorTok{=}\NormalTok{ calcular\_matriz\_confusion(y\_real, y\_pred)}
    
    \ControlFlowTok{if}\NormalTok{ TP }\OperatorTok{+}\NormalTok{ FN }\OperatorTok{==} \DecValTok{0}\NormalTok{:}
        \ControlFlowTok{return} \FloatTok{0.0}
    
\NormalTok{    recall }\OperatorTok{=}\NormalTok{ TP }\OperatorTok{/}\NormalTok{ (TP }\OperatorTok{+}\NormalTok{ FN)}
    \ControlFlowTok{return}\NormalTok{ recall}

\CommentTok{\# Calcular recall}
\NormalTok{recall }\OperatorTok{=}\NormalTok{ calcular\_recall(y\_test, y\_pred)}

\BuiltInTok{print}\NormalTok{(}\SpecialStringTok{f"Recall: }\SpecialCharTok{\{}\NormalTok{recall}\SpecialCharTok{:.4f\}}\SpecialStringTok{ (}\SpecialCharTok{\{}\NormalTok{recall}\OperatorTok{*}\DecValTok{100}\SpecialCharTok{:.2f\}}\SpecialStringTok{\%)"}\NormalTok{)}
\BuiltInTok{print}\NormalTok{(}\SpecialStringTok{f"}\CharTok{\textbackslash{}n}\SpecialStringTok{Interpretación:"}\NormalTok{)}
\BuiltInTok{print}\NormalTok{(}\SpecialStringTok{f"De los estudiantes que realmente fueron admitidos,"}\NormalTok{)}
\BuiltInTok{print}\NormalTok{(}\SpecialStringTok{f"el modelo detectó correctamente al }\SpecialCharTok{\{}\NormalTok{recall}\OperatorTok{*}\DecValTok{100}\SpecialCharTok{:.1f\}}\SpecialStringTok{\%."}\NormalTok{)}
\BuiltInTok{print}\NormalTok{(}\SpecialStringTok{f"}\CharTok{\textbackslash{}n}\SpecialStringTok{Falsos negativos: }\SpecialCharTok{\{}\NormalTok{FN}\SpecialCharTok{\}}\SpecialStringTok{ estudiantes"}\NormalTok{)}
\BuiltInTok{print}\NormalTok{(}\StringTok{"(No fueron identificados como admitidos cuando sí lo fueron)"}\NormalTok{)}
\end{Highlighting}
\end{Shaded}

\subsection{Ejemplo económico: Detección de
fraude}\label{ejemplo-econuxf3mico-detecciuxf3n-de-fraude}

\begin{Shaded}
\begin{Highlighting}[]
\BuiltInTok{print}\NormalTok{(}\StringTok{"}\CharTok{\textbackslash{}n\textbackslash{}n}\StringTok{EJEMPLO ECONÓMICO: DETECCIÓN DE FRAUDE"}\NormalTok{)}
\BuiltInTok{print}\NormalTok{(}\StringTok{"="} \OperatorTok{*} \DecValTok{70}\NormalTok{)}

\CommentTok{\# Simular transacciones}
\NormalTok{n\_trans }\OperatorTok{=} \DecValTok{10000}
\NormalTok{y\_real\_fraude }\OperatorTok{=}\NormalTok{ np.random.choice([}\DecValTok{0}\NormalTok{, }\DecValTok{1}\NormalTok{], n\_trans, p}\OperatorTok{=}\NormalTok{[}\FloatTok{0.98}\NormalTok{, }\FloatTok{0.02}\NormalTok{])  }\CommentTok{\# 2\% fraude}
\NormalTok{y\_pred\_fraude }\OperatorTok{=}\NormalTok{ np.random.choice([}\DecValTok{0}\NormalTok{, }\DecValTok{1}\NormalTok{], n\_trans, p}\OperatorTok{=}\NormalTok{[}\FloatTok{0.97}\NormalTok{, }\FloatTok{0.03}\NormalTok{])  }\CommentTok{\# Modelo detecta 3\%}

\CommentTok{\# Calcular métricas}
\NormalTok{\_, TP, TN, FP, FN }\OperatorTok{=}\NormalTok{ calcular\_matriz\_confusion(y\_real\_fraude, y\_pred\_fraude)}
\NormalTok{recall\_fraude }\OperatorTok{=}\NormalTok{ calcular\_recall(y\_real\_fraude, y\_pred\_fraude)}
\NormalTok{precision\_fraude }\OperatorTok{=}\NormalTok{ calcular\_precision(y\_real\_fraude, y\_pred\_fraude)}

\BuiltInTok{print}\NormalTok{(}\SpecialStringTok{f"}\CharTok{\textbackslash{}n}\SpecialStringTok{Resultados de }\SpecialCharTok{\{}\NormalTok{n\_trans}\SpecialCharTok{:,\}}\SpecialStringTok{ transacciones:"}\NormalTok{)}
\BuiltInTok{print}\NormalTok{(}\SpecialStringTok{f"  Fraudes reales: }\SpecialCharTok{\{}\NormalTok{TP }\OperatorTok{+}\NormalTok{ FN}\SpecialCharTok{\}}\SpecialStringTok{"}\NormalTok{)}
\BuiltInTok{print}\NormalTok{(}\SpecialStringTok{f"  Fraudes detectados: }\SpecialCharTok{\{}\NormalTok{TP}\SpecialCharTok{\}}\SpecialStringTok{ (}\SpecialCharTok{\{}\NormalTok{recall\_fraude}\OperatorTok{*}\DecValTok{100}\SpecialCharTok{:.1f\}}\SpecialStringTok{\% de los reales)"}\NormalTok{)}
\BuiltInTok{print}\NormalTok{(}\SpecialStringTok{f"  Fraudes no detectados: }\SpecialCharTok{\{}\NormalTok{FN}\SpecialCharTok{\}}\SpecialStringTok{ ¡CRÍTICO!"}\NormalTok{)}

\CommentTok{\# Calcular pérdida}
\NormalTok{monto\_promedio\_fraude }\OperatorTok{=} \DecValTok{5000}
\NormalTok{perdida\_por\_fn }\OperatorTok{=}\NormalTok{ FN }\OperatorTok{*}\NormalTok{ monto\_promedio\_fraude}

\BuiltInTok{print}\NormalTok{(}\SpecialStringTok{f"}\CharTok{\textbackslash{}n\textbackslash{}n}\SpecialStringTok{Análisis financiero:"}\NormalTok{)}
\BuiltInTok{print}\NormalTok{(}\SpecialStringTok{f"  Monto promedio por fraude: S/ }\SpecialCharTok{\{}\NormalTok{monto\_promedio\_fraude}\SpecialCharTok{:,\}}\SpecialStringTok{"}\NormalTok{)}
\BuiltInTok{print}\NormalTok{(}\SpecialStringTok{f"  Fraudes no detectados (FN): }\SpecialCharTok{\{}\NormalTok{FN}\SpecialCharTok{\}}\SpecialStringTok{"}\NormalTok{)}
\BuiltInTok{print}\NormalTok{(}\SpecialStringTok{f"  Pérdida por fraudes no detectados: S/ }\SpecialCharTok{\{}\NormalTok{perdida\_por\_fn}\SpecialCharTok{:,\}}\SpecialStringTok{"}\NormalTok{)}

\CommentTok{\# Costo de investigar falsos positivos}
\NormalTok{costo\_investigacion }\OperatorTok{=} \DecValTok{100}
\NormalTok{costo\_fp }\OperatorTok{=}\NormalTok{ FP }\OperatorTok{*}\NormalTok{ costo\_investigacion}
\BuiltInTok{print}\NormalTok{(}\SpecialStringTok{f"}\CharTok{\textbackslash{}n}\SpecialStringTok{  Alertas falsas (FP): }\SpecialCharTok{\{}\NormalTok{FP}\SpecialCharTok{\}}\SpecialStringTok{"}\NormalTok{)}
\BuiltInTok{print}\NormalTok{(}\SpecialStringTok{f"  Costo de investigar falsas alarmas: S/ }\SpecialCharTok{\{}\NormalTok{costo\_fp}\SpecialCharTok{:,\}}\SpecialStringTok{"}\NormalTok{)}

\BuiltInTok{print}\NormalTok{(}\SpecialStringTok{f"}\CharTok{\textbackslash{}n\textbackslash{}n}\SpecialStringTok{Conclusión: En detección de fraude, es MÁS IMPORTANTE el Recall"}\NormalTok{)}
\BuiltInTok{print}\NormalTok{(}\StringTok{"(No queremos perder fraudes reales, aunque tengamos falsas alarmas)"}\NormalTok{)}
\end{Highlighting}
\end{Shaded}

\section{F1-Score}\label{f1-score}

\subsection{Cálculo e
interpretación}\label{cuxe1lculo-e-interpretaciuxf3n-3}

\begin{Shaded}
\begin{Highlighting}[]
\KeywordTok{def}\NormalTok{ calcular\_f1(y\_real, y\_pred):}
    \CommentTok{"""Calcula F1{-}score"""}
\NormalTok{    precision }\OperatorTok{=}\NormalTok{ calcular\_precision(y\_real, y\_pred)}
\NormalTok{    recall }\OperatorTok{=}\NormalTok{ calcular\_recall(y\_real, y\_pred)}
    
    \ControlFlowTok{if}\NormalTok{ precision }\OperatorTok{+}\NormalTok{ recall }\OperatorTok{==} \DecValTok{0}\NormalTok{:}
        \ControlFlowTok{return} \FloatTok{0.0}
    
\NormalTok{    f1 }\OperatorTok{=} \DecValTok{2} \OperatorTok{*}\NormalTok{ (precision }\OperatorTok{*}\NormalTok{ recall) }\OperatorTok{/}\NormalTok{ (precision }\OperatorTok{+}\NormalTok{ recall)}
    \ControlFlowTok{return}\NormalTok{ f1}

\CommentTok{\# Calcular F1}
\NormalTok{f1 }\OperatorTok{=}\NormalTok{ calcular\_f1(y\_test, y\_pred)}

\BuiltInTok{print}\NormalTok{(}\SpecialStringTok{f"Métricas del modelo:"}\NormalTok{)}
\BuiltInTok{print}\NormalTok{(}\SpecialStringTok{f"  Precision: }\SpecialCharTok{\{}\NormalTok{precision}\SpecialCharTok{:.4f\}}\SpecialStringTok{"}\NormalTok{)}
\BuiltInTok{print}\NormalTok{(}\SpecialStringTok{f"  Recall: }\SpecialCharTok{\{}\NormalTok{recall}\SpecialCharTok{:.4f\}}\SpecialStringTok{"}\NormalTok{)}
\BuiltInTok{print}\NormalTok{(}\SpecialStringTok{f"  F1{-}Score: }\SpecialCharTok{\{}\NormalTok{f1}\SpecialCharTok{:.4f\}}\SpecialStringTok{"}\NormalTok{)}

\BuiltInTok{print}\NormalTok{(}\SpecialStringTok{f"}\CharTok{\textbackslash{}n\textbackslash{}n}\SpecialStringTok{Interpretación del F1{-}Score:"}\NormalTok{)}
\BuiltInTok{print}\NormalTok{(}\StringTok{"El F1{-}Score es la media armónica de Precision y Recall."}\NormalTok{)}
\BuiltInTok{print}\NormalTok{(}\SpecialStringTok{f"Un F1 de }\SpecialCharTok{\{}\NormalTok{f1}\SpecialCharTok{:.2f\}}\SpecialStringTok{ indica un }\SpecialCharTok{\{}\StringTok{\textquotesingle{}buen\textquotesingle{}} \ControlFlowTok{if}\NormalTok{ f1 }\OperatorTok{\textgreater{}} \FloatTok{0.7} \ControlFlowTok{else} \StringTok{\textquotesingle{}regular\textquotesingle{}} \ControlFlowTok{if}\NormalTok{ f1 }\OperatorTok{\textgreater{}} \FloatTok{0.5} \ControlFlowTok{else} \StringTok{\textquotesingle{}bajo\textquotesingle{}}\SpecialCharTok{\}}\SpecialStringTok{ balance entre ambas métricas."}\NormalTok{)}
\end{Highlighting}
\end{Shaded}

\subsection{Trade-off entre Precision y
Recall}\label{trade-off-entre-precision-y-recall}

\begin{Shaded}
\begin{Highlighting}[]
\CommentTok{\# Demostrar trade{-}off}
\NormalTok{umbrales }\OperatorTok{=}\NormalTok{ [}\FloatTok{0.1}\NormalTok{, }\FloatTok{0.3}\NormalTok{, }\FloatTok{0.5}\NormalTok{, }\FloatTok{0.7}\NormalTok{, }\FloatTok{0.9}\NormalTok{]}

\BuiltInTok{print}\NormalTok{(}\StringTok{"}\CharTok{\textbackslash{}n\textbackslash{}n}\StringTok{TRADE{-}OFF: PRECISION VS RECALL"}\NormalTok{)}
\BuiltInTok{print}\NormalTok{(}\StringTok{"="} \OperatorTok{*} \DecValTok{70}\NormalTok{)}
\BuiltInTok{print}\NormalTok{(}\SpecialStringTok{f"}\SpecialCharTok{\{}\StringTok{\textquotesingle{}Umbral\textquotesingle{}}\SpecialCharTok{:\textless{}10\}}\SpecialStringTok{ }\SpecialCharTok{\{}\StringTok{\textquotesingle{}Precision\textquotesingle{}}\SpecialCharTok{:\textless{}12\}}\SpecialStringTok{ }\SpecialCharTok{\{}\StringTok{\textquotesingle{}Recall\textquotesingle{}}\SpecialCharTok{:\textless{}12\}}\SpecialStringTok{ }\SpecialCharTok{\{}\StringTok{\textquotesingle{}F1{-}Score\textquotesingle{}}\SpecialCharTok{:\textless{}12\}}\SpecialStringTok{"}\NormalTok{)}
\BuiltInTok{print}\NormalTok{(}\StringTok{"{-}"} \OperatorTok{*} \DecValTok{70}\NormalTok{)}

\ControlFlowTok{for}\NormalTok{ umbral }\KeywordTok{in}\NormalTok{ umbrales:}
\NormalTok{    y\_pred\_umbral }\OperatorTok{=}\NormalTok{ (probabilidades[:, }\DecValTok{1}\NormalTok{] }\OperatorTok{\textgreater{}}\NormalTok{ umbral).astype(}\BuiltInTok{int}\NormalTok{)}
    
\NormalTok{    prec }\OperatorTok{=}\NormalTok{ calcular\_precision(y\_test, y\_pred\_umbral)}
\NormalTok{    rec }\OperatorTok{=}\NormalTok{ calcular\_recall(y\_test, y\_pred\_umbral)}
\NormalTok{    f1\_temp }\OperatorTok{=}\NormalTok{ calcular\_f1(y\_test, y\_pred\_umbral)}
    
    \BuiltInTok{print}\NormalTok{(}\SpecialStringTok{f"}\SpecialCharTok{\{}\NormalTok{umbral}\SpecialCharTok{:\textless{}10.1f\}}\SpecialStringTok{ }\SpecialCharTok{\{}\NormalTok{prec}\SpecialCharTok{:\textless{}12.4f\}}\SpecialStringTok{ }\SpecialCharTok{\{}\NormalTok{rec}\SpecialCharTok{:\textless{}12.4f\}}\SpecialStringTok{ }\SpecialCharTok{\{}\NormalTok{f1\_temp}\SpecialCharTok{:\textless{}12.4f\}}\SpecialStringTok{"}\NormalTok{)}

\BuiltInTok{print}\NormalTok{(}\StringTok{"}\CharTok{\textbackslash{}n}\StringTok{Observación:"}\NormalTok{)}
\BuiltInTok{print}\NormalTok{(}\StringTok{"{-} Umbral bajo → Mayor Recall (detecta más positivos) pero menor Precision"}\NormalTok{)}
\BuiltInTok{print}\NormalTok{(}\StringTok{"{-} Umbral alto → Mayor Precision (menos falsas alarmas) pero menor Recall"}\NormalTok{)}
\end{Highlighting}
\end{Shaded}

\section{Curva Precision-Recall}\label{curva-precision-recall}

\subsection{Generar curva
Precision-Recall}\label{generar-curva-precision-recall}

\begin{Shaded}
\begin{Highlighting}[]
\KeywordTok{def}\NormalTok{ calcular\_curva\_precision\_recall(y\_real, probabilidades):}
    \CommentTok{"""}
\CommentTok{    Calcula curva Precision{-}Recall para diferentes umbrales}
\CommentTok{    """}
    \CommentTok{\# Obtener umbrales únicos}
\NormalTok{    umbrales }\OperatorTok{=}\NormalTok{ np.sort(np.unique(probabilidades))}
    
\NormalTok{    precision\_list }\OperatorTok{=}\NormalTok{ []}
\NormalTok{    recall\_list }\OperatorTok{=}\NormalTok{ []}
    
    \ControlFlowTok{for}\NormalTok{ umbral }\KeywordTok{in}\NormalTok{ umbrales:}
\NormalTok{        y\_pred }\OperatorTok{=}\NormalTok{ (probabilidades }\OperatorTok{\textgreater{}=}\NormalTok{ umbral).astype(}\BuiltInTok{int}\NormalTok{)}
        
\NormalTok{        prec }\OperatorTok{=}\NormalTok{ calcular\_precision(y\_real, y\_pred)}
\NormalTok{        rec }\OperatorTok{=}\NormalTok{ calcular\_recall(y\_real, y\_pred)}
        
\NormalTok{        precision\_list.append(prec)}
\NormalTok{        recall\_list.append(rec)}
    
    \ControlFlowTok{return}\NormalTok{ np.array(recall\_list), np.array(precision\_list), umbrales}

\CommentTok{\# Calcular curva}
\NormalTok{recall\_curva, precision\_curva, umbrales\_pr }\OperatorTok{=}\NormalTok{ calcular\_curva\_precision\_recall(}
\NormalTok{    y\_test, probabilidades[:, }\DecValTok{1}\NormalTok{]}
\NormalTok{)}

\CommentTok{\# También usar sklearn}
\NormalTok{precision\_sklearn, recall\_sklearn, \_ }\OperatorTok{=}\NormalTok{ precision\_recall\_curve(}
\NormalTok{    y\_test, probabilidades[:, }\DecValTok{1}\NormalTok{]}
\NormalTok{)}
\end{Highlighting}
\end{Shaded}

\subsection{Visualizar curva
Precision-Recall}\label{visualizar-curva-precision-recall}

\begin{Shaded}
\begin{Highlighting}[]
\CommentTok{\# Crear gráfico}
\NormalTok{fig, ax }\OperatorTok{=}\NormalTok{ plt.subplots(figsize}\OperatorTok{=}\NormalTok{(}\DecValTok{10}\NormalTok{, }\DecValTok{8}\NormalTok{))}

\CommentTok{\# Curva propia}
\NormalTok{ax.plot(recall\_curva, precision\_curva, }\StringTok{\textquotesingle{}b{-}\textquotesingle{}}\NormalTok{, linewidth}\OperatorTok{=}\DecValTok{2}\NormalTok{, }
\NormalTok{        label}\OperatorTok{=}\StringTok{\textquotesingle{}Curva PR (implementación propia)\textquotesingle{}}\NormalTok{)}

\CommentTok{\# Curva sklearn}
\NormalTok{ax.plot(recall\_sklearn, precision\_sklearn, }\StringTok{\textquotesingle{}r{-}{-}\textquotesingle{}}\NormalTok{, linewidth}\OperatorTok{=}\DecValTok{2}\NormalTok{, alpha}\OperatorTok{=}\FloatTok{0.7}\NormalTok{,}
\NormalTok{        label}\OperatorTok{=}\StringTok{\textquotesingle{}Curva PR (sklearn)\textquotesingle{}}\NormalTok{)}

\CommentTok{\# Línea de referencia (clasificador aleatorio)}
\NormalTok{proporcion\_positivos }\OperatorTok{=}\NormalTok{ y\_test.mean()}
\NormalTok{ax.axhline(y}\OperatorTok{=}\NormalTok{proporcion\_positivos, color}\OperatorTok{=}\StringTok{\textquotesingle{}gray\textquotesingle{}}\NormalTok{, linestyle}\OperatorTok{=}\StringTok{\textquotesingle{}{-}{-}\textquotesingle{}}\NormalTok{, }
\NormalTok{           label}\OperatorTok{=}\SpecialStringTok{f\textquotesingle{}Clasificador aleatorio (}\SpecialCharTok{\{}\NormalTok{proporcion\_positivos}\SpecialCharTok{:.2f\}}\SpecialStringTok{)\textquotesingle{}}\NormalTok{)}

\CommentTok{\# Punto del modelo actual (umbral 0.5)}
\NormalTok{ax.plot(recall, precision, }\StringTok{\textquotesingle{}go\textquotesingle{}}\NormalTok{, markersize}\OperatorTok{=}\DecValTok{12}\NormalTok{, }
\NormalTok{        label}\OperatorTok{=}\SpecialStringTok{f\textquotesingle{}Umbral 0.5 (P=}\SpecialCharTok{\{}\NormalTok{precision}\SpecialCharTok{:.2f\}}\SpecialStringTok{, R=}\SpecialCharTok{\{}\NormalTok{recall}\SpecialCharTok{:.2f\}}\SpecialStringTok{)\textquotesingle{}}\NormalTok{)}

\NormalTok{ax.set\_xlabel(}\StringTok{\textquotesingle{}Recall (Sensibilidad)\textquotesingle{}}\NormalTok{, fontsize}\OperatorTok{=}\DecValTok{12}\NormalTok{)}
\NormalTok{ax.set\_ylabel(}\StringTok{\textquotesingle{}Precision (Precisión)\textquotesingle{}}\NormalTok{, fontsize}\OperatorTok{=}\DecValTok{12}\NormalTok{)}
\NormalTok{ax.set\_title(}\StringTok{\textquotesingle{}Curva Precision{-}Recall\textquotesingle{}}\NormalTok{, fontsize}\OperatorTok{=}\DecValTok{14}\NormalTok{, fontweight}\OperatorTok{=}\StringTok{\textquotesingle{}bold\textquotesingle{}}\NormalTok{)}
\NormalTok{ax.legend(loc}\OperatorTok{=}\StringTok{\textquotesingle{}best\textquotesingle{}}\NormalTok{)}
\NormalTok{ax.grid(}\VariableTok{True}\NormalTok{, alpha}\OperatorTok{=}\FloatTok{0.3}\NormalTok{)}
\NormalTok{ax.set\_xlim([}\DecValTok{0}\NormalTok{, }\FloatTok{1.05}\NormalTok{])}
\NormalTok{ax.set\_ylim([}\DecValTok{0}\NormalTok{, }\FloatTok{1.05}\NormalTok{])}

\NormalTok{plt.tight\_layout()}
\NormalTok{plt.savefig(}\StringTok{\textquotesingle{}curva\_precision\_recall.png\textquotesingle{}}\NormalTok{, dpi}\OperatorTok{=}\DecValTok{300}\NormalTok{, bbox\_inches}\OperatorTok{=}\StringTok{\textquotesingle{}tight\textquotesingle{}}\NormalTok{)}
\BuiltInTok{print}\NormalTok{(}\StringTok{"Curva Precision{-}Recall guardada como \textquotesingle{}curva\_precision\_recall.png\textquotesingle{}"}\NormalTok{)}
\end{Highlighting}
\end{Shaded}

\section{Curva ROC y AUC}\label{curva-roc-y-auc}

\subsection{Concepto de ROC}\label{concepto-de-roc}

\begin{Shaded}
\begin{Highlighting}[]
\BuiltInTok{print}\NormalTok{(}\StringTok{"""}
\StringTok{CURVA ROC (Receiver Operating Characteristic)}

\StringTok{La curva ROC grafica:}
\StringTok{{-} Eje X: False Positive Rate (FPR) = FP / (FP + TN)}
\StringTok{{-} Eje Y: True Positive Rate (TPR) = TP / (TP + FN)  [mismo que Recall]}

\StringTok{Interpretación:}
\StringTok{{-} Diagonal (línea punteada): Clasificador aleatorio (AUC = 0.5)}
\StringTok{{-} Curva sobre diagonal: Mejor que aleatorio}
\StringTok{{-} Más cerca de esquina superior izquierda: Mejor modelo}
\StringTok{{-} AUC (Area Under Curve): Métrica resumen}
\StringTok{  * AUC = 1.0: Clasificador perfecto}
\StringTok{  * AUC = 0.5: Clasificador aleatorio}
\StringTok{  * AUC \textless{} 0.5: Peor que aleatorio}
\StringTok{"""}\NormalTok{)}
\end{Highlighting}
\end{Shaded}

\subsection{Calcular TPR y FPR}\label{calcular-tpr-y-fpr}

\begin{Shaded}
\begin{Highlighting}[]
\KeywordTok{def}\NormalTok{ calcular\_tpr(y\_real, y\_pred):}
    \CommentTok{"""Calcula True Positive Rate (TPR) = Recall"""}
    \ControlFlowTok{return}\NormalTok{ calcular\_recall(y\_real, y\_pred)}

\KeywordTok{def}\NormalTok{ calcular\_fpr(y\_real, y\_pred):}
    \CommentTok{"""Calcula False Positive Rate (FPR)"""}
\NormalTok{    \_, TP, TN, FP, FN }\OperatorTok{=}\NormalTok{ calcular\_matriz\_confusion(y\_real, y\_pred)}
    
    \ControlFlowTok{if}\NormalTok{ FP }\OperatorTok{+}\NormalTok{ TN }\OperatorTok{==} \DecValTok{0}\NormalTok{:}
        \ControlFlowTok{return} \FloatTok{0.0}
    
\NormalTok{    fpr }\OperatorTok{=}\NormalTok{ FP }\OperatorTok{/}\NormalTok{ (FP }\OperatorTok{+}\NormalTok{ TN)}
    \ControlFlowTok{return}\NormalTok{ fpr}

\KeywordTok{def}\NormalTok{ calcular\_curva\_roc(y\_real, probabilidades):}
    \CommentTok{"""}
\CommentTok{    Calcula curva ROC para diferentes umbrales}
\CommentTok{    """}
\NormalTok{    umbrales }\OperatorTok{=}\NormalTok{ np.sort(np.unique(probabilidades))}
    
\NormalTok{    tpr\_list }\OperatorTok{=}\NormalTok{ []}
\NormalTok{    fpr\_list }\OperatorTok{=}\NormalTok{ []}
    
    \ControlFlowTok{for}\NormalTok{ umbral }\KeywordTok{in}\NormalTok{ umbrales:}
\NormalTok{        y\_pred }\OperatorTok{=}\NormalTok{ (probabilidades }\OperatorTok{\textgreater{}=}\NormalTok{ umbral).astype(}\BuiltInTok{int}\NormalTok{)}
        
\NormalTok{        tpr }\OperatorTok{=}\NormalTok{ calcular\_tpr(y\_real, y\_pred)}
\NormalTok{        fpr }\OperatorTok{=}\NormalTok{ calcular\_fpr(y\_real, y\_pred)}
        
\NormalTok{        tpr\_list.append(tpr)}
\NormalTok{        fpr\_list.append(fpr)}
    
    \ControlFlowTok{return}\NormalTok{ np.array(fpr\_list), np.array(tpr\_list), umbrales}

\CommentTok{\# Calcular curva ROC}
\NormalTok{fpr\_curva, tpr\_curva, umbrales\_roc }\OperatorTok{=}\NormalTok{ calcular\_curva\_roc(}
\NormalTok{    y\_test, probabilidades[:, }\DecValTok{1}\NormalTok{]}
\NormalTok{)}

\CommentTok{\# También usar sklearn}
\NormalTok{fpr\_sklearn, tpr\_sklearn, \_ }\OperatorTok{=}\NormalTok{ roc\_curve(y\_test, probabilidades[:, }\DecValTok{1}\NormalTok{])}

\CommentTok{\# Calcular AUC}
\NormalTok{roc\_auc }\OperatorTok{=}\NormalTok{ auc(fpr\_sklearn, tpr\_sklearn)}

\BuiltInTok{print}\NormalTok{(}\SpecialStringTok{f"AUC (Area Under Curve): }\SpecialCharTok{\{}\NormalTok{roc\_auc}\SpecialCharTok{:.4f\}}\SpecialStringTok{"}\NormalTok{)}
\end{Highlighting}
\end{Shaded}

\subsection{Visualizar curva ROC}\label{visualizar-curva-roc}

\begin{Shaded}
\begin{Highlighting}[]
\CommentTok{\# Crear gráfico}
\NormalTok{fig, ax }\OperatorTok{=}\NormalTok{ plt.subplots(figsize}\OperatorTok{=}\NormalTok{(}\DecValTok{10}\NormalTok{, }\DecValTok{8}\NormalTok{))}

\CommentTok{\# Curva ROC}
\NormalTok{ax.plot(fpr\_sklearn, tpr\_sklearn, }\StringTok{\textquotesingle{}b{-}\textquotesingle{}}\NormalTok{, linewidth}\OperatorTok{=}\DecValTok{2}\NormalTok{,}
\NormalTok{        label}\OperatorTok{=}\SpecialStringTok{f\textquotesingle{}Curva ROC (AUC = }\SpecialCharTok{\{}\NormalTok{roc\_auc}\SpecialCharTok{:.3f\}}\SpecialStringTok{)\textquotesingle{}}\NormalTok{)}

\CommentTok{\# Línea diagonal (clasificador aleatorio)}
\NormalTok{ax.plot([}\DecValTok{0}\NormalTok{, }\DecValTok{1}\NormalTok{], [}\DecValTok{0}\NormalTok{, }\DecValTok{1}\NormalTok{], }\StringTok{\textquotesingle{}r{-}{-}\textquotesingle{}}\NormalTok{, linewidth}\OperatorTok{=}\DecValTok{2}\NormalTok{, label}\OperatorTok{=}\StringTok{\textquotesingle{}Clasificador aleatorio (AUC = 0.5)\textquotesingle{}}\NormalTok{)}

\CommentTok{\# Punto del modelo actual}
\NormalTok{fpr\_actual }\OperatorTok{=}\NormalTok{ calcular\_fpr(y\_test, y\_pred)}
\NormalTok{tpr\_actual }\OperatorTok{=}\NormalTok{ calcular\_tpr(y\_test, y\_pred)}
\NormalTok{ax.plot(fpr\_actual, tpr\_actual, }\StringTok{\textquotesingle{}go\textquotesingle{}}\NormalTok{, markersize}\OperatorTok{=}\DecValTok{12}\NormalTok{,}
\NormalTok{        label}\OperatorTok{=}\SpecialStringTok{f\textquotesingle{}Umbral 0.5 (FPR=}\SpecialCharTok{\{}\NormalTok{fpr\_actual}\SpecialCharTok{:.2f\}}\SpecialStringTok{, TPR=}\SpecialCharTok{\{}\NormalTok{tpr\_actual}\SpecialCharTok{:.2f\}}\SpecialStringTok{)\textquotesingle{}}\NormalTok{)}

\NormalTok{ax.set\_xlabel(}\StringTok{\textquotesingle{}False Positive Rate (FPR)\textquotesingle{}}\NormalTok{, fontsize}\OperatorTok{=}\DecValTok{12}\NormalTok{)}
\NormalTok{ax.set\_ylabel(}\StringTok{\textquotesingle{}True Positive Rate (TPR)\textquotesingle{}}\NormalTok{, fontsize}\OperatorTok{=}\DecValTok{12}\NormalTok{)}
\NormalTok{ax.set\_title(}\StringTok{\textquotesingle{}Curva ROC\textquotesingle{}}\NormalTok{, fontsize}\OperatorTok{=}\DecValTok{14}\NormalTok{, fontweight}\OperatorTok{=}\StringTok{\textquotesingle{}bold\textquotesingle{}}\NormalTok{)}
\NormalTok{ax.legend(loc}\OperatorTok{=}\StringTok{\textquotesingle{}lower right\textquotesingle{}}\NormalTok{)}
\NormalTok{ax.grid(}\VariableTok{True}\NormalTok{, alpha}\OperatorTok{=}\FloatTok{0.3}\NormalTok{)}
\NormalTok{ax.set\_xlim([}\DecValTok{0}\NormalTok{, }\DecValTok{1}\NormalTok{])}
\NormalTok{ax.set\_ylim([}\DecValTok{0}\NormalTok{, }\DecValTok{1}\NormalTok{])}

\CommentTok{\# Área sombreada bajo la curva}
\NormalTok{ax.fill\_between(fpr\_sklearn, tpr\_sklearn, alpha}\OperatorTok{=}\FloatTok{0.2}\NormalTok{, label}\OperatorTok{=}\StringTok{\textquotesingle{}AUC\textquotesingle{}}\NormalTok{)}

\NormalTok{plt.tight\_layout()}
\NormalTok{plt.savefig(}\StringTok{\textquotesingle{}curva\_roc.png\textquotesingle{}}\NormalTok{, dpi}\OperatorTok{=}\DecValTok{300}\NormalTok{, bbox\_inches}\OperatorTok{=}\StringTok{\textquotesingle{}tight\textquotesingle{}}\NormalTok{)}
\BuiltInTok{print}\NormalTok{(}\StringTok{"Curva ROC guardada como \textquotesingle{}curva\_roc.png\textquotesingle{}"}\NormalTok{)}
\end{Highlighting}
\end{Shaded}

\subsection{Interpretación del AUC}\label{interpretaciuxf3n-del-auc}

\begin{Shaded}
\begin{Highlighting}[]
\KeywordTok{def}\NormalTok{ interpretar\_auc(auc\_score):}
    \CommentTok{"""Interpreta el valor de AUC"""}
    \ControlFlowTok{if}\NormalTok{ auc\_score }\OperatorTok{\textgreater{}=} \FloatTok{0.9}\NormalTok{:}
        \ControlFlowTok{return} \StringTok{"Excelente"}
    \ControlFlowTok{elif}\NormalTok{ auc\_score }\OperatorTok{\textgreater{}=} \FloatTok{0.8}\NormalTok{:}
        \ControlFlowTok{return} \StringTok{"Muy bueno"}
    \ControlFlowTok{elif}\NormalTok{ auc\_score }\OperatorTok{\textgreater{}=} \FloatTok{0.7}\NormalTok{:}
        \ControlFlowTok{return} \StringTok{"Bueno"}
    \ControlFlowTok{elif}\NormalTok{ auc\_score }\OperatorTok{\textgreater{}=} \FloatTok{0.6}\NormalTok{:}
        \ControlFlowTok{return} \StringTok{"Regular"}
    \ControlFlowTok{else}\NormalTok{:}
        \ControlFlowTok{return} \StringTok{"Pobre"}

\NormalTok{interpretacion }\OperatorTok{=}\NormalTok{ interpretar\_auc(roc\_auc)}

\BuiltInTok{print}\NormalTok{(}\SpecialStringTok{f"}\CharTok{\textbackslash{}n\textbackslash{}n}\SpecialStringTok{INTERPRETACIÓN DEL MODELO"}\NormalTok{)}
\BuiltInTok{print}\NormalTok{(}\StringTok{"="} \OperatorTok{*} \DecValTok{70}\NormalTok{)}
\BuiltInTok{print}\NormalTok{(}\SpecialStringTok{f"AUC: }\SpecialCharTok{\{}\NormalTok{roc\_auc}\SpecialCharTok{:.4f\}}\SpecialStringTok{ {-} }\SpecialCharTok{\{}\NormalTok{interpretacion}\SpecialCharTok{\}}\SpecialStringTok{"}\NormalTok{)}
\BuiltInTok{print}\NormalTok{(}\SpecialStringTok{f"}\CharTok{\textbackslash{}n}\SpecialStringTok{Significado:"}\NormalTok{)}
\BuiltInTok{print}\NormalTok{(}\SpecialStringTok{f"Si tomamos un caso positivo y uno negativo al azar,"}\NormalTok{)}
\BuiltInTok{print}\NormalTok{(}\SpecialStringTok{f"el modelo asignará una probabilidad mayor al caso positivo"}\NormalTok{)}
\BuiltInTok{print}\NormalTok{(}\SpecialStringTok{f"el }\SpecialCharTok{\{}\NormalTok{roc\_auc}\OperatorTok{*}\DecValTok{100}\SpecialCharTok{:.1f\}}\SpecialStringTok{\% de las veces."}\NormalTok{)}
\end{Highlighting}
\end{Shaded}

\section{Aplicaciones económicas}\label{aplicaciones-econuxf3micas}

\subsection{Aplicación 1: Predicción de incumplimiento
crediticio}\label{aplicaciuxf3n-1-predicciuxf3n-de-incumplimiento-crediticio}

\begin{Shaded}
\begin{Highlighting}[]
\BuiltInTok{print}\NormalTok{(}\StringTok{"}\CharTok{\textbackslash{}n\textbackslash{}n}\StringTok{APLICACIÓN: PREDICCIÓN DE INCUMPLIMIENTO CREDITICIO"}\NormalTok{)}
\BuiltInTok{print}\NormalTok{(}\StringTok{"="} \OperatorTok{*} \DecValTok{70}\NormalTok{)}

\CommentTok{\# Generar datos simulados}
\NormalTok{np.random.seed(}\DecValTok{42}\NormalTok{)}
\NormalTok{n\_clientes }\OperatorTok{=} \DecValTok{1000}

\NormalTok{datos\_credito }\OperatorTok{=}\NormalTok{ pd.DataFrame(\{}
    \StringTok{\textquotesingle{}edad\textquotesingle{}}\NormalTok{: np.random.randint(}\DecValTok{18}\NormalTok{, }\DecValTok{70}\NormalTok{, n\_clientes),}
    \StringTok{\textquotesingle{}ingreso\_mensual\textquotesingle{}}\NormalTok{: np.random.uniform(}\DecValTok{1000}\NormalTok{, }\DecValTok{10000}\NormalTok{, n\_clientes),}
    \StringTok{\textquotesingle{}monto\_credito\textquotesingle{}}\NormalTok{: np.random.uniform(}\DecValTok{5000}\NormalTok{, }\DecValTok{100000}\NormalTok{, n\_clientes),}
    \StringTok{\textquotesingle{}historial\_pagos\textquotesingle{}}\NormalTok{: np.random.randint(}\DecValTok{1}\NormalTok{, }\DecValTok{10}\NormalTok{, n\_clientes),}
    \StringTok{\textquotesingle{}num\_creditos\_activos\textquotesingle{}}\NormalTok{: np.random.randint(}\DecValTok{0}\NormalTok{, }\DecValTok{5}\NormalTok{, n\_clientes),}
    \StringTok{\textquotesingle{}ratio\_deuda\_ingreso\textquotesingle{}}\NormalTok{: np.random.uniform(}\FloatTok{0.1}\NormalTok{, }\FloatTok{0.8}\NormalTok{, n\_clientes)}
\NormalTok{\})}

\CommentTok{\# Variable objetivo: incumplimiento (1 = sí, 0 = no)}
\NormalTok{score\_riesgo }\OperatorTok{=}\NormalTok{ (}
    \OperatorTok{{-}}\NormalTok{datos\_credito[}\StringTok{\textquotesingle{}ingreso\_mensual\textquotesingle{}}\NormalTok{] }\OperatorTok{*} \FloatTok{0.0001} \OperatorTok{+}
\NormalTok{    datos\_credito[}\StringTok{\textquotesingle{}monto\_credito\textquotesingle{}}\NormalTok{] }\OperatorTok{*} \FloatTok{0.00001} \OperatorTok{+}
\NormalTok{    datos\_credito[}\StringTok{\textquotesingle{}ratio\_deuda\_ingreso\textquotesingle{}}\NormalTok{] }\OperatorTok{*} \DecValTok{2} \OperatorTok{+}
    \OperatorTok{{-}}\NormalTok{datos\_credito[}\StringTok{\textquotesingle{}historial\_pagos\textquotesingle{}}\NormalTok{] }\OperatorTok{*} \FloatTok{0.15} \OperatorTok{+}
\NormalTok{    datos\_credito[}\StringTok{\textquotesingle{}num\_creditos\_activos\textquotesingle{}}\NormalTok{] }\OperatorTok{*} \FloatTok{0.2} \OperatorTok{+}
\NormalTok{    np.random.normal(}\DecValTok{0}\NormalTok{, }\FloatTok{0.3}\NormalTok{, n\_clientes)}
\NormalTok{)}

\NormalTok{datos\_credito[}\StringTok{\textquotesingle{}incumplimiento\textquotesingle{}}\NormalTok{] }\OperatorTok{=}\NormalTok{ (score\_riesgo }\OperatorTok{\textgreater{}} \FloatTok{0.5}\NormalTok{).astype(}\BuiltInTok{int}\NormalTok{)}

\BuiltInTok{print}\NormalTok{(}\SpecialStringTok{f"}\CharTok{\textbackslash{}n}\SpecialStringTok{Dataset de créditos:"}\NormalTok{)}
\BuiltInTok{print}\NormalTok{(}\SpecialStringTok{f"  Total clientes: }\SpecialCharTok{\{}\NormalTok{n\_clientes}\SpecialCharTok{\}}\SpecialStringTok{"}\NormalTok{)}
\BuiltInTok{print}\NormalTok{(}\SpecialStringTok{f"  Incumplimientos: }\SpecialCharTok{\{}\NormalTok{datos\_credito[}\StringTok{\textquotesingle{}incumplimiento\textquotesingle{}}\NormalTok{]}\SpecialCharTok{.}\BuiltInTok{sum}\NormalTok{()}\SpecialCharTok{\}}\SpecialStringTok{ (}\SpecialCharTok{\{}\NormalTok{datos\_credito[}\StringTok{\textquotesingle{}incumplimiento\textquotesingle{}}\NormalTok{]}\SpecialCharTok{.}\NormalTok{mean()}\OperatorTok{*}\DecValTok{100}\SpecialCharTok{:.1f\}}\SpecialStringTok{\%)"}\NormalTok{)}

\CommentTok{\# Dividir datos}
\NormalTok{X\_credito }\OperatorTok{=}\NormalTok{ datos\_credito.drop(}\StringTok{\textquotesingle{}incumplimiento\textquotesingle{}}\NormalTok{, axis}\OperatorTok{=}\DecValTok{1}\NormalTok{)}
\NormalTok{y\_credito }\OperatorTok{=}\NormalTok{ datos\_credito[}\StringTok{\textquotesingle{}incumplimiento\textquotesingle{}}\NormalTok{]}

\NormalTok{X\_train\_c, X\_test\_c, y\_train\_c, y\_test\_c }\OperatorTok{=}\NormalTok{ train\_test\_split(}
\NormalTok{    X\_credito, y\_credito, test\_size}\OperatorTok{=}\FloatTok{0.3}\NormalTok{, random\_state}\OperatorTok{=}\DecValTok{42}
\NormalTok{)}

\CommentTok{\# Entrenar modelo}
\NormalTok{modelo\_credito }\OperatorTok{=}\NormalTok{ LogisticRegression(random\_state}\OperatorTok{=}\DecValTok{42}\NormalTok{, max\_iter}\OperatorTok{=}\DecValTok{1000}\NormalTok{)}
\NormalTok{modelo\_credito.fit(X\_train\_c, y\_train\_c)}

\CommentTok{\# Predecir}
\NormalTok{y\_pred\_c }\OperatorTok{=}\NormalTok{ modelo\_credito.predict(X\_test\_c)}
\NormalTok{prob\_incumplimiento }\OperatorTok{=}\NormalTok{ modelo\_credito.predict\_proba(X\_test\_c)}

\CommentTok{\# Calcular métricas}
\NormalTok{cm\_credito }\OperatorTok{=}\NormalTok{ confusion\_matrix(y\_test\_c, y\_pred\_c)}
\NormalTok{\_, TP\_c, TN\_c, FP\_c, FN\_c }\OperatorTok{=}\NormalTok{ calcular\_matriz\_confusion(y\_test\_c, y\_pred\_c)}

\NormalTok{accuracy\_c }\OperatorTok{=}\NormalTok{ calcular\_accuracy(y\_test\_c, y\_pred\_c)}
\NormalTok{precision\_c }\OperatorTok{=}\NormalTok{ calcular\_precision(y\_test\_c, y\_pred\_c)}
\NormalTok{recall\_c }\OperatorTok{=}\NormalTok{ calcular\_recall(y\_test\_c, y\_pred\_c)}
\NormalTok{f1\_c }\OperatorTok{=}\NormalTok{ calcular\_f1(y\_test\_c, y\_pred\_c)}

\CommentTok{\# Curva ROC}
\NormalTok{fpr\_c, tpr\_c, \_ }\OperatorTok{=}\NormalTok{ roc\_curve(y\_test\_c, prob\_incumplimiento[:, }\DecValTok{1}\NormalTok{])}
\NormalTok{auc\_c }\OperatorTok{=}\NormalTok{ auc(fpr\_c, tpr\_c)}

\BuiltInTok{print}\NormalTok{(}\SpecialStringTok{f"}\CharTok{\textbackslash{}n\textbackslash{}n}\SpecialStringTok{RESULTADOS DEL MODELO"}\NormalTok{)}
\BuiltInTok{print}\NormalTok{(}\StringTok{"="} \OperatorTok{*} \DecValTok{70}\NormalTok{)}
\BuiltInTok{print}\NormalTok{(}\SpecialStringTok{f"}\CharTok{\textbackslash{}n}\SpecialStringTok{Matriz de Confusión:"}\NormalTok{)}
\BuiltInTok{print}\NormalTok{(}\SpecialStringTok{f"                    Predicción"}\NormalTok{)}
\BuiltInTok{print}\NormalTok{(}\SpecialStringTok{f"              Pagará  Incumplirá"}\NormalTok{)}
\BuiltInTok{print}\NormalTok{(}\SpecialStringTok{f"Real  Pagará    }\SpecialCharTok{\{}\NormalTok{TN\_c}\SpecialCharTok{:4d\}}\SpecialStringTok{      }\SpecialCharTok{\{}\NormalTok{FP\_c}\SpecialCharTok{:4d\}}\SpecialStringTok{"}\NormalTok{)}
\BuiltInTok{print}\NormalTok{(}\SpecialStringTok{f"   Incumplirá   }\SpecialCharTok{\{}\NormalTok{FN\_c}\SpecialCharTok{:4d\}}\SpecialStringTok{      }\SpecialCharTok{\{}\NormalTok{TP\_c}\SpecialCharTok{:4d\}}\SpecialStringTok{"}\NormalTok{)}

\BuiltInTok{print}\NormalTok{(}\SpecialStringTok{f"}\CharTok{\textbackslash{}n\textbackslash{}n}\SpecialStringTok{Métricas:"}\NormalTok{)}
\BuiltInTok{print}\NormalTok{(}\SpecialStringTok{f"  Accuracy:  }\SpecialCharTok{\{}\NormalTok{accuracy\_c}\SpecialCharTok{:.4f\}}\SpecialStringTok{ (}\SpecialCharTok{\{}\NormalTok{accuracy\_c}\OperatorTok{*}\DecValTok{100}\SpecialCharTok{:.1f\}}\SpecialStringTok{\%)"}\NormalTok{)}
\BuiltInTok{print}\NormalTok{(}\SpecialStringTok{f"  Precision: }\SpecialCharTok{\{}\NormalTok{precision\_c}\SpecialCharTok{:.4f\}}\SpecialStringTok{ (}\SpecialCharTok{\{}\NormalTok{precision\_c}\OperatorTok{*}\DecValTok{100}\SpecialCharTok{:.1f\}}\SpecialStringTok{\%)"}\NormalTok{)}
\BuiltInTok{print}\NormalTok{(}\SpecialStringTok{f"  Recall:    }\SpecialCharTok{\{}\NormalTok{recall\_c}\SpecialCharTok{:.4f\}}\SpecialStringTok{ (}\SpecialCharTok{\{}\NormalTok{recall\_c}\OperatorTok{*}\DecValTok{100}\SpecialCharTok{:.1f\}}\SpecialStringTok{\%)"}\NormalTok{)}
\BuiltInTok{print}\NormalTok{(}\SpecialStringTok{f"  F1{-}Score:  }\SpecialCharTok{\{}\NormalTok{f1\_c}\SpecialCharTok{:.4f\}}\SpecialStringTok{"}\NormalTok{)}
\BuiltInTok{print}\NormalTok{(}\SpecialStringTok{f"  AUC:       }\SpecialCharTok{\{}\NormalTok{auc\_c}\SpecialCharTok{:.4f\}}\SpecialStringTok{"}\NormalTok{)}

\CommentTok{\# Análisis de costos}
\NormalTok{monto\_promedio }\OperatorTok{=}\NormalTok{ datos\_credito[}\StringTok{\textquotesingle{}monto\_credito\textquotesingle{}}\NormalTok{].mean()}
\NormalTok{tasa\_recuperacion }\OperatorTok{=} \FloatTok{0.30}  \CommentTok{\# Se recupera 30\% en promedio}

\NormalTok{costo\_fn }\OperatorTok{=}\NormalTok{ FN\_c }\OperatorTok{*}\NormalTok{ monto\_promedio }\OperatorTok{*}\NormalTok{ (}\DecValTok{1} \OperatorTok{{-}}\NormalTok{ tasa\_recuperacion)  }\CommentTok{\# No detectar incumplimiento}
\NormalTok{costo\_fp }\OperatorTok{=}\NormalTok{ FP\_c }\OperatorTok{*} \DecValTok{500}  \CommentTok{\# Costo de revisar caso que no es incumplimiento}

\BuiltInTok{print}\NormalTok{(}\SpecialStringTok{f"}\CharTok{\textbackslash{}n\textbackslash{}n}\SpecialStringTok{ANÁLISIS DE COSTOS"}\NormalTok{)}
\BuiltInTok{print}\NormalTok{(}\StringTok{"="} \OperatorTok{*} \DecValTok{70}\NormalTok{)}
\BuiltInTok{print}\NormalTok{(}\SpecialStringTok{f"  Monto promedio por crédito: S/ }\SpecialCharTok{\{}\NormalTok{monto\_promedio}\SpecialCharTok{:,.2f\}}\SpecialStringTok{"}\NormalTok{)}
\BuiltInTok{print}\NormalTok{(}\SpecialStringTok{f"}\CharTok{\textbackslash{}n}\SpecialStringTok{  Falsos Negativos (FN): }\SpecialCharTok{\{}\NormalTok{FN\_c}\SpecialCharTok{\}}\SpecialStringTok{"}\NormalTok{)}
\BuiltInTok{print}\NormalTok{(}\SpecialStringTok{f"    → Incumplimientos no detectados"}\NormalTok{)}
\BuiltInTok{print}\NormalTok{(}\SpecialStringTok{f"    → Pérdida estimada: S/ }\SpecialCharTok{\{}\NormalTok{costo\_fn}\SpecialCharTok{:,.2f\}}\SpecialStringTok{"}\NormalTok{)}
\BuiltInTok{print}\NormalTok{(}\SpecialStringTok{f"}\CharTok{\textbackslash{}n}\SpecialStringTok{  Falsos Positivos (FP): }\SpecialCharTok{\{}\NormalTok{FP\_c}\SpecialCharTok{\}}\SpecialStringTok{"}\NormalTok{)}
\BuiltInTok{print}\NormalTok{(}\SpecialStringTok{f"    → Clientes buenos rechazados o revisados"}\NormalTok{)}
\BuiltInTok{print}\NormalTok{(}\SpecialStringTok{f"    → Costo de revisión: S/ }\SpecialCharTok{\{}\NormalTok{costo\_fp}\SpecialCharTok{:,.2f\}}\SpecialStringTok{"}\NormalTok{)}
\BuiltInTok{print}\NormalTok{(}\SpecialStringTok{f"}\CharTok{\textbackslash{}n}\SpecialStringTok{  Costo total: S/ }\SpecialCharTok{\{}\NormalTok{costo\_fn }\OperatorTok{+}\NormalTok{ costo\_fp}\SpecialCharTok{:,.2f\}}\SpecialStringTok{"}\NormalTok{)}

\BuiltInTok{print}\NormalTok{(}\SpecialStringTok{f"}\CharTok{\textbackslash{}n\textbackslash{}n}\SpecialStringTok{RECOMENDACIÓN:"}\NormalTok{)}
\ControlFlowTok{if}\NormalTok{ recall\_c }\OperatorTok{\textless{}} \FloatTok{0.7}\NormalTok{:}
    \BuiltInTok{print}\NormalTok{(}\StringTok{"⚠ RECALL BAJO: Estamos perdiendo muchos incumplimientos."}\NormalTok{)}
    \BuiltInTok{print}\NormalTok{(}\StringTok{"  Considerar reducir umbral para detectar más casos de riesgo."}\NormalTok{)}
\ControlFlowTok{if}\NormalTok{ precision\_c }\OperatorTok{\textless{}} \FloatTok{0.5}\NormalTok{:}
    \BuiltInTok{print}\NormalTok{(}\StringTok{"⚠ PRECISION BAJA: Muchas falsas alarmas."}\NormalTok{)}
    \BuiltInTok{print}\NormalTok{(}\StringTok{"  Considerar aumentar umbral o mejorar features."}\NormalTok{)}
\end{Highlighting}
\end{Shaded}

\subsection{Aplicación 2: Clasificación de empresas por riesgo de
quiebra}\label{aplicaciuxf3n-2-clasificaciuxf3n-de-empresas-por-riesgo-de-quiebra}

\begin{Shaded}
\begin{Highlighting}[]
\BuiltInTok{print}\NormalTok{(}\StringTok{"}\CharTok{\textbackslash{}n\textbackslash{}n\textbackslash{}n}\StringTok{APLICACIÓN: CLASIFICACIÓN DE RIESGO DE QUIEBRA EMPRESARIAL"}\NormalTok{)}
\BuiltInTok{print}\NormalTok{(}\StringTok{"="} \OperatorTok{*} \DecValTok{70}\NormalTok{)}

\CommentTok{\# Generar datos de empresas}
\NormalTok{np.random.seed(}\DecValTok{42}\NormalTok{)}
\NormalTok{n\_empresas }\OperatorTok{=} \DecValTok{500}

\NormalTok{datos\_empresas }\OperatorTok{=}\NormalTok{ pd.DataFrame(\{}
    \StringTok{\textquotesingle{}activos\_totales\textquotesingle{}}\NormalTok{: np.random.uniform(}\DecValTok{100000}\NormalTok{, }\DecValTok{10000000}\NormalTok{, n\_empresas),}
    \StringTok{\textquotesingle{}pasivos\_totales\textquotesingle{}}\NormalTok{: np.random.uniform(}\DecValTok{50000}\NormalTok{, }\DecValTok{8000000}\NormalTok{, n\_empresas),}
    \StringTok{\textquotesingle{}ventas\_anuales\textquotesingle{}}\NormalTok{: np.random.uniform(}\DecValTok{200000}\NormalTok{, }\DecValTok{15000000}\NormalTok{, n\_empresas),}
    \StringTok{\textquotesingle{}utilidad\_neta\textquotesingle{}}\NormalTok{: np.random.uniform(}\OperatorTok{{-}}\DecValTok{500000}\NormalTok{, }\DecValTok{2000000}\NormalTok{, n\_empresas),}
    \StringTok{\textquotesingle{}flujo\_caja\textquotesingle{}}\NormalTok{: np.random.uniform(}\OperatorTok{{-}}\DecValTok{200000}\NormalTok{, }\DecValTok{1000000}\NormalTok{, n\_empresas),}
    \StringTok{\textquotesingle{}años\_operacion\textquotesingle{}}\NormalTok{: np.random.randint(}\DecValTok{1}\NormalTok{, }\DecValTok{30}\NormalTok{, n\_empresas)}
\NormalTok{\})}

\CommentTok{\# Calcular ratios financieros}
\NormalTok{datos\_empresas[}\StringTok{\textquotesingle{}ratio\_liquidez\textquotesingle{}}\NormalTok{] }\OperatorTok{=}\NormalTok{ datos\_empresas[}\StringTok{\textquotesingle{}activos\_totales\textquotesingle{}}\NormalTok{] }\OperatorTok{/}\NormalTok{ datos\_empresas[}\StringTok{\textquotesingle{}pasivos\_totales\textquotesingle{}}\NormalTok{]}
\NormalTok{datos\_empresas[}\StringTok{\textquotesingle{}margen\_utilidad\textquotesingle{}}\NormalTok{] }\OperatorTok{=}\NormalTok{ datos\_empresas[}\StringTok{\textquotesingle{}utilidad\_neta\textquotesingle{}}\NormalTok{] }\OperatorTok{/}\NormalTok{ datos\_empresas[}\StringTok{\textquotesingle{}ventas\_anuales\textquotesingle{}}\NormalTok{]}
\NormalTok{datos\_empresas[}\StringTok{\textquotesingle{}ratio\_endeudamiento\textquotesingle{}}\NormalTok{] }\OperatorTok{=}\NormalTok{ datos\_empresas[}\StringTok{\textquotesingle{}pasivos\_totales\textquotesingle{}}\NormalTok{] }\OperatorTok{/}\NormalTok{ datos\_empresas[}\StringTok{\textquotesingle{}activos\_totales\textquotesingle{}}\NormalTok{]}

\CommentTok{\# Variable objetivo: quiebra (1 = sí, 0 = no)}
\NormalTok{score\_quiebra }\OperatorTok{=}\NormalTok{ (}
    \OperatorTok{{-}}\NormalTok{datos\_empresas[}\StringTok{\textquotesingle{}ratio\_liquidez\textquotesingle{}}\NormalTok{] }\OperatorTok{*} \FloatTok{0.5} \OperatorTok{+}
    \OperatorTok{{-}}\NormalTok{datos\_empresas[}\StringTok{\textquotesingle{}margen\_utilidad\textquotesingle{}}\NormalTok{] }\OperatorTok{*} \DecValTok{2} \OperatorTok{+}
\NormalTok{    datos\_empresas[}\StringTok{\textquotesingle{}ratio\_endeudamiento\textquotesingle{}}\NormalTok{] }\OperatorTok{*} \DecValTok{3} \OperatorTok{+}
    \OperatorTok{{-}}\NormalTok{datos\_empresas[}\StringTok{\textquotesingle{}flujo\_caja\textquotesingle{}}\NormalTok{] }\OperatorTok{*} \FloatTok{0.000001} \OperatorTok{+}
    \OperatorTok{{-}}\NormalTok{datos\_empresas[}\StringTok{\textquotesingle{}años\_operacion\textquotesingle{}}\NormalTok{] }\OperatorTok{*} \FloatTok{0.05} \OperatorTok{+}
\NormalTok{    np.random.normal(}\DecValTok{0}\NormalTok{, }\FloatTok{0.5}\NormalTok{, n\_empresas)}
\NormalTok{)}

\NormalTok{datos\_empresas[}\StringTok{\textquotesingle{}quiebra\textquotesingle{}}\NormalTok{] }\OperatorTok{=}\NormalTok{ (score\_quiebra }\OperatorTok{\textgreater{}} \FloatTok{1.0}\NormalTok{).astype(}\BuiltInTok{int}\NormalTok{)}

\BuiltInTok{print}\NormalTok{(}\SpecialStringTok{f"}\CharTok{\textbackslash{}n}\SpecialStringTok{Dataset de empresas:"}\NormalTok{)}
\BuiltInTok{print}\NormalTok{(}\SpecialStringTok{f"  Total empresas: }\SpecialCharTok{\{}\NormalTok{n\_empresas}\SpecialCharTok{\}}\SpecialStringTok{"}\NormalTok{)}
\BuiltInTok{print}\NormalTok{(}\SpecialStringTok{f"  Quiebras: }\SpecialCharTok{\{}\NormalTok{datos\_empresas[}\StringTok{\textquotesingle{}quiebra\textquotesingle{}}\NormalTok{]}\SpecialCharTok{.}\BuiltInTok{sum}\NormalTok{()}\SpecialCharTok{\}}\SpecialStringTok{ (}\SpecialCharTok{\{}\NormalTok{datos\_empresas[}\StringTok{\textquotesingle{}quiebra\textquotesingle{}}\NormalTok{]}\SpecialCharTok{.}\NormalTok{mean()}\OperatorTok{*}\DecValTok{100}\SpecialCharTok{:.1f\}}\SpecialStringTok{\%)"}\NormalTok{)}

\CommentTok{\# Seleccionar features}
\NormalTok{features\_empresas }\OperatorTok{=}\NormalTok{ [}\StringTok{\textquotesingle{}ratio\_liquidez\textquotesingle{}}\NormalTok{, }\StringTok{\textquotesingle{}margen\_utilidad\textquotesingle{}}\NormalTok{, }\StringTok{\textquotesingle{}ratio\_endeudamiento\textquotesingle{}}\NormalTok{,}
                     \StringTok{\textquotesingle{}flujo\_caja\textquotesingle{}}\NormalTok{, }\StringTok{\textquotesingle{}años\_operacion\textquotesingle{}}\NormalTok{]}

\NormalTok{X\_empresas }\OperatorTok{=}\NormalTok{ datos\_empresas[features\_empresas]}
\NormalTok{y\_empresas }\OperatorTok{=}\NormalTok{ datos\_empresas[}\StringTok{\textquotesingle{}quiebra\textquotesingle{}}\NormalTok{]}

\CommentTok{\# Dividir datos}
\NormalTok{X\_train\_e, X\_test\_e, y\_train\_e, y\_test\_e }\OperatorTok{=}\NormalTok{ train\_test\_split(}
\NormalTok{    X\_empresas, y\_empresas, test\_size}\OperatorTok{=}\FloatTok{0.3}\NormalTok{, random\_state}\OperatorTok{=}\DecValTok{42}
\NormalTok{)}

\CommentTok{\# Entrenar modelo}
\NormalTok{modelo\_quiebra }\OperatorTok{=}\NormalTok{ LogisticRegression(random\_state}\OperatorTok{=}\DecValTok{42}\NormalTok{, max\_iter}\OperatorTok{=}\DecValTok{1000}\NormalTok{)}
\NormalTok{modelo\_quiebra.fit(X\_train\_e, y\_train\_e)}

\CommentTok{\# Evaluar}
\NormalTok{y\_pred\_e }\OperatorTok{=}\NormalTok{ modelo\_quiebra.predict(X\_test\_e)}
\NormalTok{prob\_quiebra }\OperatorTok{=}\NormalTok{ modelo\_quiebra.predict\_proba(X\_test\_e)}

\CommentTok{\# Métricas}
\NormalTok{accuracy\_e }\OperatorTok{=}\NormalTok{ calcular\_accuracy(y\_test\_e, y\_pred\_e)}
\NormalTok{precision\_e }\OperatorTok{=}\NormalTok{ calcular\_precision(y\_test\_e, y\_pred\_e)}
\NormalTok{recall\_e }\OperatorTok{=}\NormalTok{ calcular\_recall(y\_test\_e, y\_pred\_e)}
\NormalTok{f1\_e }\OperatorTok{=}\NormalTok{ calcular\_f1(y\_test\_e, y\_pred\_e)}

\NormalTok{fpr\_e, tpr\_e, \_ }\OperatorTok{=}\NormalTok{ roc\_curve(y\_test\_e, prob\_quiebra[:, }\DecValTok{1}\NormalTok{])}
\NormalTok{auc\_e }\OperatorTok{=}\NormalTok{ auc(fpr\_e, tpr\_e)}

\BuiltInTok{print}\NormalTok{(}\SpecialStringTok{f"}\CharTok{\textbackslash{}n\textbackslash{}n}\SpecialStringTok{MÉTRICAS DEL MODELO"}\NormalTok{)}
\BuiltInTok{print}\NormalTok{(}\StringTok{"="} \OperatorTok{*} \DecValTok{70}\NormalTok{)}
\BuiltInTok{print}\NormalTok{(}\SpecialStringTok{f"  Accuracy:  }\SpecialCharTok{\{}\NormalTok{accuracy\_e}\SpecialCharTok{:.4f\}}\SpecialStringTok{"}\NormalTok{)}
\BuiltInTok{print}\NormalTok{(}\SpecialStringTok{f"  Precision: }\SpecialCharTok{\{}\NormalTok{precision\_e}\SpecialCharTok{:.4f\}}\SpecialStringTok{"}\NormalTok{)}
\BuiltInTok{print}\NormalTok{(}\SpecialStringTok{f"  Recall:    }\SpecialCharTok{\{}\NormalTok{recall\_e}\SpecialCharTok{:.4f\}}\SpecialStringTok{"}\NormalTok{)}
\BuiltInTok{print}\NormalTok{(}\SpecialStringTok{f"  F1{-}Score:  }\SpecialCharTok{\{}\NormalTok{f1\_e}\SpecialCharTok{:.4f\}}\SpecialStringTok{"}\NormalTok{)}
\BuiltInTok{print}\NormalTok{(}\SpecialStringTok{f"  AUC:       }\SpecialCharTok{\{}\NormalTok{auc\_e}\SpecialCharTok{:.4f\}}\SpecialStringTok{"}\NormalTok{)}

\CommentTok{\# Importancia de variables}
\BuiltInTok{print}\NormalTok{(}\SpecialStringTok{f"}\CharTok{\textbackslash{}n\textbackslash{}n}\SpecialStringTok{IMPORTANCIA DE VARIABLES (Coeficientes)"}\NormalTok{)}
\BuiltInTok{print}\NormalTok{(}\StringTok{"="} \OperatorTok{*} \DecValTok{70}\NormalTok{)}
\NormalTok{coeficientes }\OperatorTok{=}\NormalTok{ pd.DataFrame(\{}
    \StringTok{\textquotesingle{}Variable\textquotesingle{}}\NormalTok{: features\_empresas,}
    \StringTok{\textquotesingle{}Coeficiente\textquotesingle{}}\NormalTok{: modelo\_quiebra.coef\_[}\DecValTok{0}\NormalTok{]}
\NormalTok{\}).sort\_values(}\StringTok{\textquotesingle{}Coeficiente\textquotesingle{}}\NormalTok{, key}\OperatorTok{=}\BuiltInTok{abs}\NormalTok{, ascending}\OperatorTok{=}\VariableTok{False}\NormalTok{)}

\ControlFlowTok{for}\NormalTok{ \_, row }\KeywordTok{in}\NormalTok{ coeficientes.iterrows():}
\NormalTok{    signo }\OperatorTok{=} \StringTok{"↑ Mayor riesgo"} \ControlFlowTok{if}\NormalTok{ row[}\StringTok{\textquotesingle{}Coeficiente\textquotesingle{}}\NormalTok{] }\OperatorTok{\textgreater{}} \DecValTok{0} \ControlFlowTok{else} \StringTok{"↓ Menor riesgo"}
    \BuiltInTok{print}\NormalTok{(}\SpecialStringTok{f"  }\SpecialCharTok{\{}\NormalTok{row[}\StringTok{\textquotesingle{}Variable\textquotesingle{}}\NormalTok{]}\SpecialCharTok{:\textless{}25\}}\SpecialStringTok{ }\SpecialCharTok{\{}\NormalTok{row[}\StringTok{\textquotesingle{}Coeficiente\textquotesingle{}}\NormalTok{]}\SpecialCharTok{:\textgreater{}8.4f\}}\SpecialStringTok{  }\SpecialCharTok{\{}\NormalTok{signo}\SpecialCharTok{\}}\SpecialStringTok{"}\NormalTok{)}

\BuiltInTok{print}\NormalTok{(}\SpecialStringTok{f"}\CharTok{\textbackslash{}n\textbackslash{}n}\SpecialStringTok{APLICACIÓN PRÁCTICA:"}\NormalTok{)}
\BuiltInTok{print}\NormalTok{(}\StringTok{"Este modelo puede usarse para:"}\NormalTok{)}
\BuiltInTok{print}\NormalTok{(}\StringTok{"  1. Sistema de alerta temprana de quiebra"}\NormalTok{)}
\BuiltInTok{print}\NormalTok{(}\StringTok{"  2. Decisiones de inversión o préstamos corporativos"}\NormalTok{)}
\BuiltInTok{print}\NormalTok{(}\StringTok{"  3. Priorización de auditorías o intervenciones"}\NormalTok{)}
\BuiltInTok{print}\NormalTok{(}\StringTok{"  4. Análisis de cartera de inversiones"}\NormalTok{)}
\end{Highlighting}
\end{Shaded}

\section{Ejercicios prácticos}\label{ejercicios-pruxe1cticos}

\subsection{Ejercicio 1: Optimizar umbral según
costo}\label{ejercicio-1-optimizar-umbral-seguxfan-costo}

\begin{Shaded}
\begin{Highlighting}[]
\BuiltInTok{print}\NormalTok{(}\StringTok{"}\CharTok{\textbackslash{}n\textbackslash{}n\textbackslash{}n}\StringTok{EJERCICIO: OPTIMIZACIÓN DE UMBRAL SEGÚN COSTOS"}\NormalTok{)}
\BuiltInTok{print}\NormalTok{(}\StringTok{"="} \OperatorTok{*} \DecValTok{70}\NormalTok{)}

\KeywordTok{def}\NormalTok{ encontrar\_umbral\_optimo(y\_real, probabilidades, costo\_fp, costo\_fn):}
    \CommentTok{"""}
\CommentTok{    Encuentra el umbral óptimo que minimiza costos totales}
\CommentTok{    """}
\NormalTok{    umbrales }\OperatorTok{=}\NormalTok{ np.linspace(}\FloatTok{0.1}\NormalTok{, }\FloatTok{0.9}\NormalTok{, }\DecValTok{50}\NormalTok{)}
    
\NormalTok{    mejores\_resultados }\OperatorTok{=}\NormalTok{ \{}
        \StringTok{\textquotesingle{}umbral\textquotesingle{}}\NormalTok{: }\VariableTok{None}\NormalTok{,}
        \StringTok{\textquotesingle{}costo\_total\textquotesingle{}}\NormalTok{: }\BuiltInTok{float}\NormalTok{(}\StringTok{\textquotesingle{}inf\textquotesingle{}}\NormalTok{),}
        \StringTok{\textquotesingle{}precision\textquotesingle{}}\NormalTok{: }\DecValTok{0}\NormalTok{,}
        \StringTok{\textquotesingle{}recall\textquotesingle{}}\NormalTok{: }\DecValTok{0}\NormalTok{,}
        \StringTok{\textquotesingle{}fp\textquotesingle{}}\NormalTok{: }\DecValTok{0}\NormalTok{,}
        \StringTok{\textquotesingle{}fn\textquotesingle{}}\NormalTok{: }\DecValTok{0}
\NormalTok{    \}}
    
\NormalTok{    resultados }\OperatorTok{=}\NormalTok{ []}
    
    \ControlFlowTok{for}\NormalTok{ umbral }\KeywordTok{in}\NormalTok{ umbrales:}
\NormalTok{        y\_pred }\OperatorTok{=}\NormalTok{ (probabilidades }\OperatorTok{\textgreater{}=}\NormalTok{ umbral).astype(}\BuiltInTok{int}\NormalTok{)}
        
\NormalTok{        \_, TP, TN, FP, FN }\OperatorTok{=}\NormalTok{ calcular\_matriz\_confusion(y\_real, y\_pred)}
        
\NormalTok{        costo\_total }\OperatorTok{=}\NormalTok{ FP }\OperatorTok{*}\NormalTok{ costo\_fp }\OperatorTok{+}\NormalTok{ FN }\OperatorTok{*}\NormalTok{ costo\_fn}
        
\NormalTok{        prec }\OperatorTok{=}\NormalTok{ calcular\_precision(y\_real, y\_pred)}
\NormalTok{        rec }\OperatorTok{=}\NormalTok{ calcular\_recall(y\_real, y\_pred)}
        
\NormalTok{        resultados.append(\{}
            \StringTok{\textquotesingle{}umbral\textquotesingle{}}\NormalTok{: umbral,}
            \StringTok{\textquotesingle{}costo\_total\textquotesingle{}}\NormalTok{: costo\_total,}
            \StringTok{\textquotesingle{}precision\textquotesingle{}}\NormalTok{: prec,}
            \StringTok{\textquotesingle{}recall\textquotesingle{}}\NormalTok{: rec,}
            \StringTok{\textquotesingle{}fp\textquotesingle{}}\NormalTok{: FP,}
            \StringTok{\textquotesingle{}fn\textquotesingle{}}\NormalTok{: FN}
\NormalTok{        \})}
        
        \ControlFlowTok{if}\NormalTok{ costo\_total }\OperatorTok{\textless{}}\NormalTok{ mejores\_resultados[}\StringTok{\textquotesingle{}costo\_total\textquotesingle{}}\NormalTok{]:}
\NormalTok{            mejores\_resultados }\OperatorTok{=}\NormalTok{ \{}
                \StringTok{\textquotesingle{}umbral\textquotesingle{}}\NormalTok{: umbral,}
                \StringTok{\textquotesingle{}costo\_total\textquotesingle{}}\NormalTok{: costo\_total,}
                \StringTok{\textquotesingle{}precision\textquotesingle{}}\NormalTok{: prec,}
                \StringTok{\textquotesingle{}recall\textquotesingle{}}\NormalTok{: rec,}
                \StringTok{\textquotesingle{}fp\textquotesingle{}}\NormalTok{: FP,}
                \StringTok{\textquotesingle{}fn\textquotesingle{}}\NormalTok{: FN}
\NormalTok{            \}}
    
    \ControlFlowTok{return}\NormalTok{ mejores\_resultados, pd.DataFrame(resultados)}

\CommentTok{\# Aplicar al modelo de crédito}
\NormalTok{costo\_por\_fp }\OperatorTok{=} \DecValTok{500}  \CommentTok{\# Costo de revisar un caso falso positivo}
\NormalTok{costo\_por\_fn }\OperatorTok{=}\NormalTok{ monto\_promedio }\OperatorTok{*} \FloatTok{0.7}  \CommentTok{\# 70\% del monto no recuperado}

\NormalTok{resultado\_optimo, df\_resultados }\OperatorTok{=}\NormalTok{ encontrar\_umbral\_optimo(}
\NormalTok{    y\_test\_c, }
\NormalTok{    prob\_incumplimiento[:, }\DecValTok{1}\NormalTok{],}
\NormalTok{    costo\_por\_fp,}
\NormalTok{    costo\_por\_fn}
\NormalTok{)}

\BuiltInTok{print}\NormalTok{(}\SpecialStringTok{f"}\CharTok{\textbackslash{}n}\SpecialStringTok{Costos considerados:"}\NormalTok{)}
\BuiltInTok{print}\NormalTok{(}\SpecialStringTok{f"  Costo por Falso Positivo: S/ }\SpecialCharTok{\{}\NormalTok{costo\_por\_fp}\SpecialCharTok{:,.2f\}}\SpecialStringTok{"}\NormalTok{)}
\BuiltInTok{print}\NormalTok{(}\SpecialStringTok{f"  Costo por Falso Negativo: S/ }\SpecialCharTok{\{}\NormalTok{costo\_por\_fn}\SpecialCharTok{:,.2f\}}\SpecialStringTok{"}\NormalTok{)}

\BuiltInTok{print}\NormalTok{(}\SpecialStringTok{f"}\CharTok{\textbackslash{}n\textbackslash{}n}\SpecialStringTok{UMBRAL ÓPTIMO ENCONTRADO:"}\NormalTok{)}
\BuiltInTok{print}\NormalTok{(}\StringTok{"="} \OperatorTok{*} \DecValTok{70}\NormalTok{)}
\BuiltInTok{print}\NormalTok{(}\SpecialStringTok{f"  Umbral óptimo: }\SpecialCharTok{\{}\NormalTok{resultado\_optimo[}\StringTok{\textquotesingle{}umbral\textquotesingle{}}\NormalTok{]}\SpecialCharTok{:.3f\}}\SpecialStringTok{"}\NormalTok{)}
\BuiltInTok{print}\NormalTok{(}\SpecialStringTok{f"  Costo total: S/ }\SpecialCharTok{\{}\NormalTok{resultado\_optimo[}\StringTok{\textquotesingle{}costo\_total\textquotesingle{}}\NormalTok{]}\SpecialCharTok{:,.2f\}}\SpecialStringTok{"}\NormalTok{)}
\BuiltInTok{print}\NormalTok{(}\SpecialStringTok{f"  Precision: }\SpecialCharTok{\{}\NormalTok{resultado\_optimo[}\StringTok{\textquotesingle{}precision\textquotesingle{}}\NormalTok{]}\SpecialCharTok{:.4f\}}\SpecialStringTok{"}\NormalTok{)}
\BuiltInTok{print}\NormalTok{(}\SpecialStringTok{f"  Recall: }\SpecialCharTok{\{}\NormalTok{resultado\_optimo[}\StringTok{\textquotesingle{}recall\textquotesingle{}}\NormalTok{]}\SpecialCharTok{:.4f\}}\SpecialStringTok{"}\NormalTok{)}
\BuiltInTok{print}\NormalTok{(}\SpecialStringTok{f"  Falsos Positivos: }\SpecialCharTok{\{}\NormalTok{resultado\_optimo[}\StringTok{\textquotesingle{}fp\textquotesingle{}}\NormalTok{]}\SpecialCharTok{\}}\SpecialStringTok{"}\NormalTok{)}
\BuiltInTok{print}\NormalTok{(}\SpecialStringTok{f"  Falsos Negativos: }\SpecialCharTok{\{}\NormalTok{resultado\_optimo[}\StringTok{\textquotesingle{}fn\textquotesingle{}}\NormalTok{]}\SpecialCharTok{\}}\SpecialStringTok{"}\NormalTok{)}

\CommentTok{\# Comparar con umbral 0.5}
\NormalTok{\_, TP\_05, TN\_05, FP\_05, FN\_05 }\OperatorTok{=}\NormalTok{ calcular\_matriz\_confusion(y\_test\_c, y\_pred\_c)}
\NormalTok{costo\_umbral\_05 }\OperatorTok{=}\NormalTok{ FP\_05 }\OperatorTok{*}\NormalTok{ costo\_por\_fp }\OperatorTok{+}\NormalTok{ FN\_05 }\OperatorTok{*}\NormalTok{ costo\_por\_fn}

\BuiltInTok{print}\NormalTok{(}\SpecialStringTok{f"}\CharTok{\textbackslash{}n\textbackslash{}n}\SpecialStringTok{COMPARACIÓN CON UMBRAL 0.5:"}\NormalTok{)}
\BuiltInTok{print}\NormalTok{(}\SpecialStringTok{f"  Costo con umbral 0.5: S/ }\SpecialCharTok{\{}\NormalTok{costo\_umbral\_05}\SpecialCharTok{:,.2f\}}\SpecialStringTok{"}\NormalTok{)}
\BuiltInTok{print}\NormalTok{(}\SpecialStringTok{f"  Ahorro con umbral óptimo: S/ }\SpecialCharTok{\{}\NormalTok{costo\_umbral\_05 }\OperatorTok{{-}}\NormalTok{ resultado\_optimo[}\StringTok{\textquotesingle{}costo\_total\textquotesingle{}}\NormalTok{]}\SpecialCharTok{:,.2f\}}\SpecialStringTok{"}\NormalTok{)}

\CommentTok{\# Graficar}
\NormalTok{fig, (ax1, ax2) }\OperatorTok{=}\NormalTok{ plt.subplots(}\DecValTok{1}\NormalTok{, }\DecValTok{2}\NormalTok{, figsize}\OperatorTok{=}\NormalTok{(}\DecValTok{15}\NormalTok{, }\DecValTok{5}\NormalTok{))}

\CommentTok{\# Gráfico de costos}
\NormalTok{ax1.plot(df\_resultados[}\StringTok{\textquotesingle{}umbral\textquotesingle{}}\NormalTok{], df\_resultados[}\StringTok{\textquotesingle{}costo\_total\textquotesingle{}}\NormalTok{], }\StringTok{\textquotesingle{}b{-}\textquotesingle{}}\NormalTok{, linewidth}\OperatorTok{=}\DecValTok{2}\NormalTok{)}
\NormalTok{ax1.axvline(x}\OperatorTok{=}\NormalTok{resultado\_optimo[}\StringTok{\textquotesingle{}umbral\textquotesingle{}}\NormalTok{], color}\OperatorTok{=}\StringTok{\textquotesingle{}r\textquotesingle{}}\NormalTok{, linestyle}\OperatorTok{=}\StringTok{\textquotesingle{}{-}{-}\textquotesingle{}}\NormalTok{, }
\NormalTok{            label}\OperatorTok{=}\SpecialStringTok{f"Umbral óptimo (}\SpecialCharTok{\{}\NormalTok{resultado\_optimo[}\StringTok{\textquotesingle{}umbral\textquotesingle{}}\NormalTok{]}\SpecialCharTok{:.3f\}}\SpecialStringTok{)"}\NormalTok{)}
\NormalTok{ax1.axvline(x}\OperatorTok{=}\FloatTok{0.5}\NormalTok{, color}\OperatorTok{=}\StringTok{\textquotesingle{}g\textquotesingle{}}\NormalTok{, linestyle}\OperatorTok{=}\StringTok{\textquotesingle{}{-}{-}\textquotesingle{}}\NormalTok{, alpha}\OperatorTok{=}\FloatTok{0.5}\NormalTok{, label}\OperatorTok{=}\StringTok{\textquotesingle{}Umbral 0.5\textquotesingle{}}\NormalTok{)}
\NormalTok{ax1.set\_xlabel(}\StringTok{\textquotesingle{}Umbral\textquotesingle{}}\NormalTok{)}
\NormalTok{ax1.set\_ylabel(}\StringTok{\textquotesingle{}Costo Total (S/)\textquotesingle{}}\NormalTok{)}
\NormalTok{ax1.set\_title(}\StringTok{\textquotesingle{}Costo Total vs Umbral\textquotesingle{}}\NormalTok{)}
\NormalTok{ax1.legend()}
\NormalTok{ax1.grid(}\VariableTok{True}\NormalTok{, alpha}\OperatorTok{=}\FloatTok{0.3}\NormalTok{)}

\CommentTok{\# Gráfico de Precision vs Recall}
\NormalTok{ax2.plot(df\_resultados[}\StringTok{\textquotesingle{}umbral\textquotesingle{}}\NormalTok{], df\_resultados[}\StringTok{\textquotesingle{}precision\textquotesingle{}}\NormalTok{], }\StringTok{\textquotesingle{}b{-}\textquotesingle{}}\NormalTok{, }
\NormalTok{         linewidth}\OperatorTok{=}\DecValTok{2}\NormalTok{, label}\OperatorTok{=}\StringTok{\textquotesingle{}Precision\textquotesingle{}}\NormalTok{)}
\NormalTok{ax2.plot(df\_resultados[}\StringTok{\textquotesingle{}umbral\textquotesingle{}}\NormalTok{], df\_resultados[}\StringTok{\textquotesingle{}recall\textquotesingle{}}\NormalTok{], }\StringTok{\textquotesingle{}r{-}\textquotesingle{}}\NormalTok{, }
\NormalTok{         linewidth}\OperatorTok{=}\DecValTok{2}\NormalTok{, label}\OperatorTok{=}\StringTok{\textquotesingle{}Recall\textquotesingle{}}\NormalTok{)}
\NormalTok{ax2.axvline(x}\OperatorTok{=}\NormalTok{resultado\_optimo[}\StringTok{\textquotesingle{}umbral\textquotesingle{}}\NormalTok{], color}\OperatorTok{=}\StringTok{\textquotesingle{}g\textquotesingle{}}\NormalTok{, linestyle}\OperatorTok{=}\StringTok{\textquotesingle{}{-}{-}\textquotesingle{}}\NormalTok{, alpha}\OperatorTok{=}\FloatTok{0.5}\NormalTok{,}
\NormalTok{            label}\OperatorTok{=}\StringTok{\textquotesingle{}Umbral óptimo\textquotesingle{}}\NormalTok{)}
\NormalTok{ax2.set\_xlabel(}\StringTok{\textquotesingle{}Umbral\textquotesingle{}}\NormalTok{)}
\NormalTok{ax2.set\_ylabel(}\StringTok{\textquotesingle{}Métrica\textquotesingle{}}\NormalTok{)}
\NormalTok{ax2.set\_title(}\StringTok{\textquotesingle{}Precision y Recall vs Umbral\textquotesingle{}}\NormalTok{)}
\NormalTok{ax2.legend()}
\NormalTok{ax2.grid(}\VariableTok{True}\NormalTok{, alpha}\OperatorTok{=}\FloatTok{0.3}\NormalTok{)}

\NormalTok{plt.tight\_layout()}
\NormalTok{plt.savefig(}\StringTok{\textquotesingle{}optimizacion\_umbral.png\textquotesingle{}}\NormalTok{, dpi}\OperatorTok{=}\DecValTok{300}\NormalTok{, bbox\_inches}\OperatorTok{=}\StringTok{\textquotesingle{}tight\textquotesingle{}}\NormalTok{)}
\BuiltInTok{print}\NormalTok{(}\StringTok{"}\CharTok{\textbackslash{}n}\StringTok{Gráfico guardado como \textquotesingle{}optimizacion\_umbral.png\textquotesingle{}"}\NormalTok{)}
\end{Highlighting}
\end{Shaded}

\section{Conclusión}\label{conclusiuxf3n}

En esta guía hemos explorado predicción y métricas de performance:

\textbf{El problema de predicción}

Predecir eventos futuros basándose en datos históricos para tomar
mejores decisiones.

\textbf{Matriz de confusión}

Herramienta fundamental que muestra TP, TN, FP, FN para evaluar
clasificadores.

\textbf{Métricas de clasificación}

\begin{itemize}
\tightlist
\item
  Accuracy: Proporción de predicciones correctas
\item
  Precision: De las predicciones positivas, cuántas son correctas
\item
  Recall: De los casos positivos reales, cuántos detectamos
\item
  F1-Score: Balance entre precision y recall
\end{itemize}

\textbf{Curvas de evaluación}

\begin{itemize}
\tightlist
\item
  Curva Precision-Recall: Para evaluar trade-off entre precision y
  recall
\item
  Curva ROC: Para evaluar capacidad discriminativa del modelo
\item
  AUC: Métrica resumen de la curva ROC
\end{itemize}

\section{Próximos pasos}\label{pruxf3ximos-pasos}

En la siguiente guía (Guía 9: Métodos de ML para Clasificación)
exploraremos:

\begin{itemize}
\tightlist
\item
  Regresión logística
\item
  K-Nearest Neighbors (KNN)
\item
  Árboles de decisión
\item
  Random Forests
\item
  Clasificación multi-clase
\end{itemize}

Estos algoritmos nos permitirán construir modelos predictivos más
sofisticados para problemas económicos complejos.

\section{Recursos adicionales}\label{recursos-adicionales}

Para profundizar en métricas de clasificación:

\begin{itemize}
\tightlist
\item
  Scikit-learn Metrics:
  scikit-learn.org/stable/modules/model\_evaluation.html
\item
  Understanding ROC Curves:
  developers.google.com/machine-learning/crash-course/classification/roc-and-auc
\item
  Precision-Recall trade-off: machinelearningmastery.com
\item
  Aplicaciones en economía: papers en credit scoring, fraud detection
\end{itemize}

\section{Publicaciones Similares}\label{publicaciones-similares}

Si te interesó este artículo, te recomendamos que explores otros blogs y
recursos relacionados que pueden ampliar tus conocimientos. Aquí te dejo
algunas sugerencias:

\begin{enumerate}
\def\labelenumi{\arabic{enumi}.}
\tightlist
\item
  \href{https://numerus-scriptum.netlify.app/python/2020-06-19-instalacion-de-anaconda/index.pdf}{\faIcon{file-pdf}}
  \href{https://numerus-scriptum.netlify.app/python/2020-06-19-instalacion-de-anaconda}{Instalacion
  De Anaconda}
\item
  \href{https://numerus-scriptum.netlify.app/python/2020-06-20-configurar-entorno-virtual-python-anaconda/index.pdf}{\faIcon{file-pdf}}
  \href{https://numerus-scriptum.netlify.app/python/2020-06-20-configurar-entorno-virtual-python-anaconda}{Configurar
  Entorno Virtual Python Anaconda}
\item
  \href{https://numerus-scriptum.netlify.app/python/2021-04-17-01-introducion-a-la-programacion-con-python/index.pdf}{\faIcon{file-pdf}}
  \href{https://numerus-scriptum.netlify.app/python/2021-04-17-01-introducion-a-la-programacion-con-python}{01
  Introducion A La Programacion Con Python}
\item
  \href{https://numerus-scriptum.netlify.app/python/2021-05-31-02-variables-expresiones-y-statements-con-python/index.pdf}{\faIcon{file-pdf}}
  \href{https://numerus-scriptum.netlify.app/python/2021-05-31-02-variables-expresiones-y-statements-con-python}{02
  Variables Expresiones Y Statements Con Python}
\item
  \href{https://numerus-scriptum.netlify.app/python/2021-06-07-03-objetos-de-python/index.pdf}{\faIcon{file-pdf}}
  \href{https://numerus-scriptum.netlify.app/python/2021-06-07-03-objetos-de-python}{03
  Objetos De Python}
\item
  \href{https://numerus-scriptum.netlify.app/python/2021-06-14-04-ejecucion-condicional-con-python/index.pdf}{\faIcon{file-pdf}}
  \href{https://numerus-scriptum.netlify.app/python/2021-06-14-04-ejecucion-condicional-con-python}{04
  Ejecucion Condicional Con Python}
\item
  \href{https://numerus-scriptum.netlify.app/python/2021-06-21-05-iteraciones-con-python/index.pdf}{\faIcon{file-pdf}}
  \href{https://numerus-scriptum.netlify.app/python/2021-06-21-05-iteraciones-con-python}{05
  Iteraciones Con Python}
\item
  \href{https://numerus-scriptum.netlify.app/python/2021-08-16-06-funciones-con-python/index.pdf}{\faIcon{file-pdf}}
  \href{https://numerus-scriptum.netlify.app/python/2021-08-16-06-funciones-con-python}{06
  Funciones Con Python}
\item
  \href{https://numerus-scriptum.netlify.app/python/2021-08-23-07-dataframes-con-python/index.pdf}{\faIcon{file-pdf}}
  \href{https://numerus-scriptum.netlify.app/python/2021-08-23-07-dataframes-con-python}{07
  Dataframes Con Python}
\item
  \href{https://numerus-scriptum.netlify.app/python/2021-11-29-08-prediccion-y-metrica-de-performance-con-python/index.pdf}{\faIcon{file-pdf}}
  \href{https://numerus-scriptum.netlify.app/python/2021-11-29-08-prediccion-y-metrica-de-performance-con-python}{08
  Prediccion Y Metrica De Performance Con Python}
\item
  \href{https://numerus-scriptum.netlify.app/python/2021-12-06-09-metodos-de-machine-learning-para-clasificacion-con-python/index.pdf}{\faIcon{file-pdf}}
  \href{https://numerus-scriptum.netlify.app/python/2021-12-06-09-metodos-de-machine-learning-para-clasificacion-con-python}{09
  Metodos De Machine Learning Para Clasificacion Con Python}
\item
  \href{https://numerus-scriptum.netlify.app/python/2021-12-13-10-metodos-de-machine-learning-para-regresion-con-python/index.pdf}{\faIcon{file-pdf}}
  \href{https://numerus-scriptum.netlify.app/python/2021-12-13-10-metodos-de-machine-learning-para-regresion-con-python}{10
  Metodos De Machine Learning Para Regresion Con Python}
\item
  \href{https://numerus-scriptum.netlify.app/python/2022-10-31-11-validacion-cruzada-y-composicion-del-modelo-con-python/index.pdf}{\faIcon{file-pdf}}
  \href{https://numerus-scriptum.netlify.app/python/2022-10-31-11-validacion-cruzada-y-composicion-del-modelo-con-python}{11
  Validacion Cruzada Y Composicion Del Modelo Con Python}
\item
  \href{https://numerus-scriptum.netlify.app/python/2025-05-10-visualizacion-de-datos-con-python/index.pdf}{\faIcon{file-pdf}}
  \href{https://numerus-scriptum.netlify.app/python/2025-05-10-visualizacion-de-datos-con-python}{Visualizacion
  De Datos Con Python}
\end{enumerate}

Esperamos que encuentres estas publicaciones igualmente interesantes y
útiles. ¡Disfruta de la lectura!






\end{document}
